% ============================================================
% Chapter 5: Group-Based Activity-Oriented Relax-and-Fix
% Thesis: Solving the FSRCPSP Using Hybrid ML Approaches
% Author: Souvik Chakraborty
% ============================================================
\chapter{Group-Based Activity-Oriented Relax-and-Fix}
\label{chap:grouprf}
The baseline R\&F procedure (Chapter~\ref{chap:baseline}) fixes
one activity per subproblem, yielding weak LP relaxations and an
index-based processing order that ignores activity similarity
(Section~\ref{sec:baseline_limits}). This chapter introduces
the group-based extension, which addresses both limitations by
embedding an unsupervised machine learning pipeline that
identifies groups of semantically similar activities and fixes
each group simultaneously in a single subproblem.

% ============================================================
% TikZ diagram: Group-Based Activity-Oriented R&F Pipeline
% (Updated to match AO_FR_Gurobi_updated.py)
%
% Include in Chapter 5 with:
%   % ============================================================
% TikZ diagram: Group-Based Activity-Oriented R&F Pipeline
% (Updated to match AO_FR_Gurobi_updated.py)
%
% Include in Chapter 5 with:
%   % ============================================================
% TikZ diagram: Group-Based Activity-Oriented R&F Pipeline
% (Updated to match AO_FR_Gurobi_updated.py)
%
% Include in Chapter 5 with:
%   \input{figures/grouprf_pipeline.tex}
%
% Required in preamble:
%   \usepackage{tikz}
%   \usetikzlibrary{shapes.geometric, arrows.meta,
%                   positioning, fit, backgrounds}
% ============================================================
\begin{figure}[p]
\centering
\vspace*{\fill}
\resizebox{\linewidth}{!}{%
\begin{tikzpicture}[
    node distance     = 0.7cm and 1.5cm,
    box/.style        = {rectangle, rounded corners=6pt,
                         draw=black!30, fill=gray!8,
                         minimum width=8cm, minimum height=0.9cm,
                         align=center, font=\small},
    boxteal/.style    = {box, fill=teal!10,   draw=teal!50},
    boxpurp/.style    = {box, fill=violet!8,  draw=violet!40},
    boxblue/.style    = {box, fill=blue!8,    draw=blue!40},
    boxamb/.style     = {rectangle, rounded corners=14pt,
                         draw=orange!50, fill=orange!8,
                         minimum width=8cm, minimum height=0.9cm,
                         align=center, font=\small},
    arr/.style        = {-{Stealth[length=5pt]}, thick,
                         draw=black!50},
    darr/.style       = {-{Stealth[length=5pt]}, draw=black!30,
                         dashed, thin},
    lbl/.style        = {font=\scriptsize, text=black!60},
]
%% ── Nodes ──────────────────────────────────────────────────────────
\node[box] (inst)
    {Project instance $\mathcal{N}=(V,E)$, $K$, $S$, $T$};
\node[boxteal, below=of inst] (ssgs)
    {\textbf{SSGS}: greedy initial schedule $S$};
%% ── ML grouping block ──────────────────────────────────────────────
\node[boxpurp, below=of ssgs] (feat)
    {\textbf{Enriched feature extraction + redundancy pruning}\\[2pt]
     \small \texttt{EnrichedFeatureExtractor}
            $\;\longrightarrow\;\tilde{X}$
            $\;\cdot\;$ drop features with $|r| \geq 0.95$};
\node[boxpurp, below=of feat] (cluster)
    {\textbf{Ward clustering per precedence level}\\[2pt]
     \small Longest-path DP $\rightarrow$ levels
            $\;\cdot\;$ Silhouette score $\rightarrow k^*$
            $\;\cdot\;$ merge adjacent groups};
\node[box, below=of cluster] (groups)
    {Groups $\mathcal{G} = [G_1, G_2, \ldots, G_m]$};
\node[lbl] (mllabel) at ([yshift=11pt]feat.north)
    {ML grouping pipeline};
\begin{scope}[on background layer]
    \node[draw=black!20, dashed, rounded corners=8pt,
          inner sep=12pt, fill=violet!3,
          fit=(feat)(cluster)(groups)(mllabel)]
          (mlbox) {};
\end{scope}
%% ── R&F loop block ─────────────────────────────────────────────────
\node[boxblue, below=2.0cm of groups] (window)
    {\textbf{Narrow slack-based window} for each $i \in G_g$\\[2pt]
     \small $w_i = \max\!\left(2,
             \left\lfloor 0.15\,(LS_i - ES_i)\right\rfloor\right)$,
             \quad centred on $S'(i)$};
\node[boxblue, below=of window] (subprob)
    {\textbf{Solve group subproblem} (SAA-FSRCPSP)\\[2pt]
     \small Binary: $G_g$
            $\;\cdot\;$ Fixed: $G_1,\ldots,G_{g-1}$
            $\;\cdot\;$ Relaxed: remainder};
\node[boxamb, below=of subprob] (decision)
    {Feasible $S'$ and improves makespan?};
\node[boxteal, below=of decision] (fix)
    {\textbf{Fix} $G_g$ at $S'$
     $\;\cdot\;$ \textbf{Update} $S = S'$
     $\;\cdot\;$ $\mathrm{no\_improve} = 0$};
\node[lbl] (rflabel) at ([yshift=11pt]window.north)
    {Group-based R\&F loop ($g = 1,\ldots,m$)};
\begin{scope}[on background layer]
    \node[draw=black!15, dashed, rounded corners=8pt,
          inner sep=10pt, fill=blue!2,
          fit=(window)(subprob)(decision)(fix)(rflabel)]
          (rfbox) {};
\end{scope}
%% ── Arrows ─────────────────────────────────────────────────────────
\draw[arr] (inst)     -- (ssgs);
\draw[arr] (ssgs)     -- (feat);
\draw[arr] (feat)     -- (cluster);
\draw[arr] (cluster)  -- (groups);
\draw[arr] (groups)   -- (window);
\draw[arr] (window)   -- (subprob);
\draw[arr] (subprob)  -- (decision);
\draw[arr] (decision) -- node[lbl, right]{yes} (fix);
%% ── No: full-slack retry / fix at S' fallback ──────────────────────
\draw[darr]
    (decision.west)
    -- ++(-2.2, 0)
    |- (window.west)
    node[lbl, align=left, pos=0.25, left]
        {No,\\full-slack\\retry or fix\\at $S'$};
%% ── Loop next group ─────────────────────────────────────────────────
\draw[darr]
    (fix.east)
    -- ++(2.2, 0)
    |- (window.east)
    node[lbl, pos=0.25, right]{next $g$};
\end{tikzpicture}
}
\caption{Pipeline of the group-based activity-oriented Relax-and-Fix
         (updated).
         The ML grouping pipeline (dashed violet box) uses enriched
         feature extraction with correlation-based redundancy pruning,
         precedence-level assignment via longest-path DP, and Ward
         hierarchical clustering with silhouette-based dynamic cut
         to form semantically coherent groups.
         The R\&F loop (dashed blue box) processes one group per
         subproblem with a two-stage window strategy: a narrow
         slack-based window ($15\%$ of slack) is tried first;
         if no solution is found, the full $[ES_i, LS_i]$ window
         is used as fallback.  If both solves fail, the group is
         fixed at the current best schedule~$S'$.
         Early stopping triggers when an integral solution is found
         or after consecutive non-improving groups.}
\label{fig:grouprf_pipeline}
\vspace*{\fill}
\end{figure}

%
% Required in preamble:
%   \usepackage{tikz}
%   \usetikzlibrary{shapes.geometric, arrows.meta,
%                   positioning, fit, backgrounds}
% ============================================================
\begin{figure}[p]
\centering
\vspace*{\fill}
\resizebox{\linewidth}{!}{%
\begin{tikzpicture}[
    node distance     = 0.7cm and 1.5cm,
    box/.style        = {rectangle, rounded corners=6pt,
                         draw=black!30, fill=gray!8,
                         minimum width=8cm, minimum height=0.9cm,
                         align=center, font=\small},
    boxteal/.style    = {box, fill=teal!10,   draw=teal!50},
    boxpurp/.style    = {box, fill=violet!8,  draw=violet!40},
    boxblue/.style    = {box, fill=blue!8,    draw=blue!40},
    boxamb/.style     = {rectangle, rounded corners=14pt,
                         draw=orange!50, fill=orange!8,
                         minimum width=8cm, minimum height=0.9cm,
                         align=center, font=\small},
    arr/.style        = {-{Stealth[length=5pt]}, thick,
                         draw=black!50},
    darr/.style       = {-{Stealth[length=5pt]}, draw=black!30,
                         dashed, thin},
    lbl/.style        = {font=\scriptsize, text=black!60},
]
%% ── Nodes ──────────────────────────────────────────────────────────
\node[box] (inst)
    {Project instance $\mathcal{N}=(V,E)$, $K$, $S$, $T$};
\node[boxteal, below=of inst] (ssgs)
    {\textbf{SSGS}: greedy initial schedule $S$};
%% ── ML grouping block ──────────────────────────────────────────────
\node[boxpurp, below=of ssgs] (feat)
    {\textbf{Enriched feature extraction + redundancy pruning}\\[2pt]
     \small \texttt{EnrichedFeatureExtractor}
            $\;\longrightarrow\;\tilde{X}$
            $\;\cdot\;$ drop features with $|r| \geq 0.95$};
\node[boxpurp, below=of feat] (cluster)
    {\textbf{Ward clustering per precedence level}\\[2pt]
     \small Longest-path DP $\rightarrow$ levels
            $\;\cdot\;$ Silhouette score $\rightarrow k^*$
            $\;\cdot\;$ merge adjacent groups};
\node[box, below=of cluster] (groups)
    {Groups $\mathcal{G} = [G_1, G_2, \ldots, G_m]$};
\node[lbl] (mllabel) at ([yshift=11pt]feat.north)
    {ML grouping pipeline};
\begin{scope}[on background layer]
    \node[draw=black!20, dashed, rounded corners=8pt,
          inner sep=12pt, fill=violet!3,
          fit=(feat)(cluster)(groups)(mllabel)]
          (mlbox) {};
\end{scope}
%% ── R&F loop block ─────────────────────────────────────────────────
\node[boxblue, below=2.0cm of groups] (window)
    {\textbf{Narrow slack-based window} for each $i \in G_g$\\[2pt]
     \small $w_i = \max\!\left(2,
             \left\lfloor 0.15\,(LS_i - ES_i)\right\rfloor\right)$,
             \quad centred on $S'(i)$};
\node[boxblue, below=of window] (subprob)
    {\textbf{Solve group subproblem} (SAA-FSRCPSP)\\[2pt]
     \small Binary: $G_g$
            $\;\cdot\;$ Fixed: $G_1,\ldots,G_{g-1}$
            $\;\cdot\;$ Relaxed: remainder};
\node[boxamb, below=of subprob] (decision)
    {Feasible $S'$ and improves makespan?};
\node[boxteal, below=of decision] (fix)
    {\textbf{Fix} $G_g$ at $S'$
     $\;\cdot\;$ \textbf{Update} $S = S'$
     $\;\cdot\;$ $\mathrm{no\_improve} = 0$};
\node[lbl] (rflabel) at ([yshift=11pt]window.north)
    {Group-based R\&F loop ($g = 1,\ldots,m$)};
\begin{scope}[on background layer]
    \node[draw=black!15, dashed, rounded corners=8pt,
          inner sep=10pt, fill=blue!2,
          fit=(window)(subprob)(decision)(fix)(rflabel)]
          (rfbox) {};
\end{scope}
%% ── Arrows ─────────────────────────────────────────────────────────
\draw[arr] (inst)     -- (ssgs);
\draw[arr] (ssgs)     -- (feat);
\draw[arr] (feat)     -- (cluster);
\draw[arr] (cluster)  -- (groups);
\draw[arr] (groups)   -- (window);
\draw[arr] (window)   -- (subprob);
\draw[arr] (subprob)  -- (decision);
\draw[arr] (decision) -- node[lbl, right]{yes} (fix);
%% ── No: full-slack retry / fix at S' fallback ──────────────────────
\draw[darr]
    (decision.west)
    -- ++(-2.2, 0)
    |- (window.west)
    node[lbl, align=left, pos=0.25, left]
        {No,\\full-slack\\retry or fix\\at $S'$};
%% ── Loop next group ─────────────────────────────────────────────────
\draw[darr]
    (fix.east)
    -- ++(2.2, 0)
    |- (window.east)
    node[lbl, pos=0.25, right]{next $g$};
\end{tikzpicture}
}
\caption{Pipeline of the group-based activity-oriented Relax-and-Fix
         (updated).
         The ML grouping pipeline (dashed violet box) uses enriched
         feature extraction with correlation-based redundancy pruning,
         precedence-level assignment via longest-path DP, and Ward
         hierarchical clustering with silhouette-based dynamic cut
         to form semantically coherent groups.
         The R\&F loop (dashed blue box) processes one group per
         subproblem with a two-stage window strategy: a narrow
         slack-based window ($15\%$ of slack) is tried first;
         if no solution is found, the full $[ES_i, LS_i]$ window
         is used as fallback.  If both solves fail, the group is
         fixed at the current best schedule~$S'$.
         Early stopping triggers when an integral solution is found
         or after consecutive non-improving groups.}
\label{fig:grouprf_pipeline}
\vspace*{\fill}
\end{figure}

%
% Required in preamble:
%   \usepackage{tikz}
%   \usetikzlibrary{shapes.geometric, arrows.meta,
%                   positioning, fit, backgrounds}
% ============================================================
\begin{figure}[p]
\centering
\vspace*{\fill}
\resizebox{\linewidth}{!}{%
\begin{tikzpicture}[
    node distance     = 0.7cm and 1.5cm,
    box/.style        = {rectangle, rounded corners=6pt,
                         draw=black!30, fill=gray!8,
                         minimum width=8cm, minimum height=0.9cm,
                         align=center, font=\small},
    boxteal/.style    = {box, fill=teal!10,   draw=teal!50},
    boxpurp/.style    = {box, fill=violet!8,  draw=violet!40},
    boxblue/.style    = {box, fill=blue!8,    draw=blue!40},
    boxamb/.style     = {rectangle, rounded corners=14pt,
                         draw=orange!50, fill=orange!8,
                         minimum width=8cm, minimum height=0.9cm,
                         align=center, font=\small},
    arr/.style        = {-{Stealth[length=5pt]}, thick,
                         draw=black!50},
    darr/.style       = {-{Stealth[length=5pt]}, draw=black!30,
                         dashed, thin},
    lbl/.style        = {font=\scriptsize, text=black!60},
]
%% ── Nodes ──────────────────────────────────────────────────────────
\node[box] (inst)
    {Project instance $\mathcal{N}=(V,E)$, $K$, $S$, $T$};
\node[boxteal, below=of inst] (ssgs)
    {\textbf{SSGS}: greedy initial schedule $S$};
%% ── ML grouping block ──────────────────────────────────────────────
\node[boxpurp, below=of ssgs] (feat)
    {\textbf{Enriched feature extraction + redundancy pruning}\\[2pt]
     \small \texttt{EnrichedFeatureExtractor}
            $\;\longrightarrow\;\tilde{X}$
            $\;\cdot\;$ drop features with $|r| \geq 0.95$};
\node[boxpurp, below=of feat] (cluster)
    {\textbf{Ward clustering per precedence level}\\[2pt]
     \small Longest-path DP $\rightarrow$ levels
            $\;\cdot\;$ Silhouette score $\rightarrow k^*$
            $\;\cdot\;$ merge adjacent groups};
\node[box, below=of cluster] (groups)
    {Groups $\mathcal{G} = [G_1, G_2, \ldots, G_m]$};
\node[lbl] (mllabel) at ([yshift=11pt]feat.north)
    {ML grouping pipeline};
\begin{scope}[on background layer]
    \node[draw=black!20, dashed, rounded corners=8pt,
          inner sep=12pt, fill=violet!3,
          fit=(feat)(cluster)(groups)(mllabel)]
          (mlbox) {};
\end{scope}
%% ── R&F loop block ─────────────────────────────────────────────────
\node[boxblue, below=2.0cm of groups] (window)
    {\textbf{Narrow slack-based window} for each $i \in G_g$\\[2pt]
     \small $w_i = \max\!\left(2,
             \left\lfloor 0.15\,(LS_i - ES_i)\right\rfloor\right)$,
             \quad centred on $S'(i)$};
\node[boxblue, below=of window] (subprob)
    {\textbf{Solve group subproblem} (SAA-FSRCPSP)\\[2pt]
     \small Binary: $G_g$
            $\;\cdot\;$ Fixed: $G_1,\ldots,G_{g-1}$
            $\;\cdot\;$ Relaxed: remainder};
\node[boxamb, below=of subprob] (decision)
    {Feasible $S'$ and improves makespan?};
\node[boxteal, below=of decision] (fix)
    {\textbf{Fix} $G_g$ at $S'$
     $\;\cdot\;$ \textbf{Update} $S = S'$
     $\;\cdot\;$ $\mathrm{no\_improve} = 0$};
\node[lbl] (rflabel) at ([yshift=11pt]window.north)
    {Group-based R\&F loop ($g = 1,\ldots,m$)};
\begin{scope}[on background layer]
    \node[draw=black!15, dashed, rounded corners=8pt,
          inner sep=10pt, fill=blue!2,
          fit=(window)(subprob)(decision)(fix)(rflabel)]
          (rfbox) {};
\end{scope}
%% ── Arrows ─────────────────────────────────────────────────────────
\draw[arr] (inst)     -- (ssgs);
\draw[arr] (ssgs)     -- (feat);
\draw[arr] (feat)     -- (cluster);
\draw[arr] (cluster)  -- (groups);
\draw[arr] (groups)   -- (window);
\draw[arr] (window)   -- (subprob);
\draw[arr] (subprob)  -- (decision);
\draw[arr] (decision) -- node[lbl, right]{yes} (fix);
%% ── No: full-slack retry / fix at S' fallback ──────────────────────
\draw[darr]
    (decision.west)
    -- ++(-2.2, 0)
    |- (window.west)
    node[lbl, align=left, pos=0.25, left]
        {No,\\full-slack\\retry or fix\\at $S'$};
%% ── Loop next group ─────────────────────────────────────────────────
\draw[darr]
    (fix.east)
    -- ++(2.2, 0)
    |- (window.east)
    node[lbl, pos=0.25, right]{next $g$};
\end{tikzpicture}
}
\caption{Pipeline of the group-based activity-oriented Relax-and-Fix
         (updated).
         The ML grouping pipeline (dashed violet box) uses enriched
         feature extraction with correlation-based redundancy pruning,
         precedence-level assignment via longest-path DP, and Ward
         hierarchical clustering with silhouette-based dynamic cut
         to form semantically coherent groups.
         The R\&F loop (dashed blue box) processes one group per
         subproblem with a two-stage window strategy: a narrow
         slack-based window ($15\%$ of slack) is tried first;
         if no solution is found, the full $[ES_i, LS_i]$ window
         is used as fallback.  If both solves fail, the group is
         fixed at the current best schedule~$S'$.
         Early stopping triggers when an integral solution is found
         or after consecutive non-improving groups.}
\label{fig:grouprf_pipeline}
\vspace*{\fill}
\end{figure}

% ------------------------------------------------------------
\section{Overview and Motivation}
\label{sec:grouprf_overview}
% ------------------------------------------------------------
The decomposition of binary variables into subsets that are fixed
iteratively is the core principle of Fix-and-Optimize~(F\&O) and
Relax-and-Fix~(R\&F) methods \parencite{pochet2006, Helber2010}. The
key design decision in any R\&F implementation is how to define
the subsets. \textcite{Helber2010} distinguished between product-,
resource-, and process-oriented decomposition strategies for the
capacitated lot-sizing problem, showing that the decomposition
structure significantly affects solution quality. For project
scheduling, \textcite{bigler2024multisite} demonstrated that a unit-oriented
decomposition, which fixes one sequencing decision per iteration,
produces strong results on the multi-site RCPSP.
% The baseline approach of \citet{mendl2024} adopts an activity-oriented
% decomposition in which each activity is its own unit.
The group-based approach proposed in this thesis extends this
activity-oriented decomposition by grouping activities according
to their feature similarity before applying the R\&F procedure.
The underlying idea is that fixing a group of related activities
simultaneously gives Gurobi more binary variables in
a single subproblem, which tends to produce tighter LP relaxations
and stronger root-node bounds \parencite{pochet2006}. At the same
time, the grouping step draws on unsupervised clustering, placing
this work within the broader trend of using learned structure to
accelerate MIP solving \parencite{bengio2021}.
Concretely, each group $G_g \subseteq V$ is formed by the
unsupervised pipeline and processed as a single subproblem,
reducing the total number of subproblems from $n+2$ (baseline)
to $m = |\mathcal{G}|$, where $m \ll n$ for typical instances.
The complete pipeline (Figure~\ref{fig:grouprf_pipeline})
consists of three components: feature extraction
(Section~\ref{sec:features}), activity grouping
(Section~\ref{sec:grouping}), and the group-based R\&F
decomposition (Section~\ref{sec:group_rf_decomp}).
% ------------------------------------------------------------
\section{Feature Extraction}
\label{sec:features}
% ------------------------------------------------------------
To enable similarity-based grouping, each activity $i \in V$ is
described by a feature vector drawn from a comprehensive candidate
pool of $18 + 8|K|$ features ($50$ for $|K|=4$). The pool is
organised into seven groups spanning five aspects of scheduling
behaviour: network position (Groups~A--B, $7$~features),
duration (Group~C, $3$~features), scheduling urgency
(Group~D, $3$~features), timing (Group~E, $4$~features), and
stochastic resource load (Groups~F--G, $1 + 8|K|$~features).
All candidate features are computed by the
\texttt{EnrichedFeatureExtractor} class. Feature engineering
for project scheduling has roots in the priority rule literature
\parencite{Kolisch1999}; the candidate pool here is
intentionally broad: it includes every plausible activity
descriptor that can be constructed from the instance data, and
a comprehensive ablation
(Chapter~\ref{chap:experiments},
Section~\ref{sec:ablation}) then determines which
features actually help the clustering
\parencite{Guyon2003}.

\subsection{Network Position Features (Group~A)}
\label{subsec:net_features}
Four features capture the structural position of each activity
within the precedence network $\mathcal{N} = (V, E)$. These
features collectively characterise an activity's position
relative to the flow of work: how far from the project start it
sits, how far from the project end, how many activities must
converge before it can begin, and how many activities it enables
upon completion.
% ----------------------------------------------------------------
\subsubsection*{Precedence Level}
\label{subsubsec:prec_level}
% ----------------------------------------------------------------
The \textit{precedence level} $\ell_i$ is defined as the length of the
longest path from the source node $0$ to activity $i$ in the precedence
network:
\begin{equation}
    \ell_i = \max_{j : (j,i) \in E} (\ell_j + 1),
    \qquad \ell_0 = 0
    \label{eq:prec_level}
\end{equation}
This quantity is computed in two stages. First, Kahn's topological sort
algorithm \parencite{Kahn1962} produces a valid linear ordering of $V$
such that for every edge $(j,i) \in E$, activity $j$ precedes activity
$i$ in the ordering. This ordering is a necessary precondition for the
dynamic programming recurrence in Equation~\eqref{eq:prec_level}, since
$\ell_i$ can only be evaluated once all predecessor values $\ell_j$ are
already known.
\begin{algorithm}[ht]
\caption{Kahn's topological sort \parencite{Kahn1962}}
\label{alg:kahn}
\begin{algorithmic}[1]
\State Compute $\mathrm{in\_deg}[v] \leftarrow
       |\{u : (u,v) \in E\}|$ for all $v \in V$
\State $Q \leftarrow \{v \in V : \mathrm{in\_deg}[v] = 0\}$
\State $L \leftarrow [\,]$
\While{$Q \neq \emptyset$}
    \State $u \leftarrow \textsc{Dequeue}(Q)$,\;
           append $u$ to $L$
    \For{each $v$ with $(u,v) \in E$}
        \State $\mathrm{in\_deg}[v] \leftarrow
               \mathrm{in\_deg}[v] - 1$
        \If{$\mathrm{in\_deg}[v] = 0$}
            \State $\textsc{Enqueue}(Q,\, v)$
        \EndIf
    \EndFor
\EndWhile
\State \Return topological order $L$
\end{algorithmic}
\end{algorithm}
Once the topological order $L = [v_1, v_2, \ldots, v_{|V|}]$
is obtained, the longest-path dynamic programme is applied by
processing nodes strictly in this order, as given in
Algorithm~\ref{alg:prec_level_dp}.

\begin{algorithm}[ht]
\caption{Longest-path DP for precedence levels}
\label{alg:prec_level_dp}
\begin{algorithmic}[1]
\State $\ell[0] \leftarrow 0$
\For{each activity $i$ in topological order $L$,
     excluding source node $0$}
    \State $\ell[i] \leftarrow
           \max\bigl\{\ell[j] + 1 : (j,\,i) \in E\bigr\}$
\EndFor
\State \Return $\ell[i]$ for all $i \in V$
\end{algorithmic}
\end{algorithm}


Structurally, this is the most important feature for grouping, since
activities sharing the same precedence level, i.e.\
$\ell_i = \ell_j$ for $i \neq j$, are guaranteed to have no
direct or transitive precedence relationship between them.
This follows directly from the longest-path DP: if a directed
path of length $p \geq 1$ existed from $i$ to $k$, then
$\ell_k \geq \ell_i + p > \ell_i$, contradicting
$\ell_i = \ell_k$. This property is a necessary condition
for safe simultaneous fixing in the MIP subproblem; fixing
two activities that stand in a precedence relationship could
yield an infeasible partial schedule. Precedence level
therefore serves as the primary partitioning criterion in
the grouping pipeline (see Section~\ref{subsec:prec_safety}).

\subsubsection*{In-Degree}
\label{subsubsec:in_degree}

The \textit{in-degree} of activity $i$ is the number of its
direct predecessors in the precedence network:

\begin{equation}
    \mathrm{in\_deg}_i = |\{j : (j,i) \in E\}|
    \label{eq:in_deg}
\end{equation}

A simple arc-counting pass computes this for all activities.
Activities with high in-degree are \textit{synchronisation
points}: they cannot start until multiple parallel chains
of work have all finished \parencite{Demeulemeester2002}. Any
delay in a single predecessor propagates directly to them,
making them natural bottlenecks. Two activities at the same
precedence level may still differ substantially in in-degree,
so this feature adds information that $\ell_i$ alone does not
capture.

\subsubsection*{Out-Degree}
\label{subsubsec:out_degree}

The \textit{out-degree} of activity $i$ is the number of
activities it directly enables upon completion:

\begin{equation}
    \mathrm{out\_deg}_i = |\{j : (i,j) \in E\}|
    \label{eq:out_deg}
\end{equation}

Activities with high out-degree act as \textit{enabling hubs}:
their completion unlocks many parallel successors
\parencite{Demeulemeester2002}. Scheduling such an activity
late has a wide downstream ripple effect on the makespan,
so the R\&F heuristic benefits from fixing them early.
% ----------------------------------------------------------------
\subsubsection*{Longest Path to Sink}
\label{subsubsec:longest_path_sink}
% ----------------------------------------------------------------
The \textit{longest path to sink} $\delta_i$ is defined as the
length of the longest path from activity $i$ to the project sink
node $n+1$:
\begin{equation}
    \delta_i = \max_{j : (i,j) \in E} (\delta_j + 1),
    \qquad \delta_{n+1} = 0
    \label{eq:longest_path_sink}
\end{equation}
This is computed by a backward dynamic programme that processes
the topological order in reverse, as given in
Algorithm~\ref{alg:backward_dp}.

\begin{algorithm}[ht]
\caption{Backward longest-path DP for depth to sink}
\label{alg:backward_dp}
\begin{algorithmic}[1]
\State $\delta[n+1] \leftarrow 0$
\For{each activity $j$ in \textbf{reverse} topological order $L$}
    \For{each successor $s$ of $j$}
        \State $\delta[j] \leftarrow
               \max\bigl(\delta[j],\; \delta[s] + 1\bigr)$
    \EndFor
\EndFor
\State \Return $\delta[i]$ for all $i \in V$
\end{algorithmic}
\end{algorithm}

Pairing $\delta_i$ with $\ell_i$ pins each activity from both
ends of the network. An activity with low $\ell_i$ and high
$\delta_i$ sits near the start of a long chain; one with high
$\ell_i$ and low $\delta_i$ is close to the project end. Two
activities at the same precedence level can still occupy very
different positions along parallel chains of different lengths,
and $\delta_i$ is what distinguishes them.

\subsubsection*{Summary}

Together, the four features provide a compact but informative
structural signature for each activity within the precedence
network. Table~\ref{tab:net_features} summarises their
definitions and scheduling interpretation.

\begin{table}[ht]
\centering
\renewcommand{\arraystretch}{1.4}
\caption{Network position features and their scheduling
         interpretation}
\label{tab:net_features}
\begin{tabular}{llp{5.5cm}}
\hline
\textbf{Feature} &
\textbf{Formula} &
\textbf{Scheduling interpretation} \\
\hline
Precedence level $\ell_i$
    & $\max_{j:(j,i)\in E}(\ell_j+1)$
    & Depth from source; primary
      partitioning criterion for safe grouping \\[4pt]
In-degree $\mathrm{in\_deg}_i$
    & $|\{j:(j,i)\in E\}|$
    & Degree of upstream convergence;
      identifies synchronisation bottlenecks \\[4pt]
Out-degree $\mathrm{out\_deg}_i$
    & $|\{j:(i,j)\in E\}|$
    & Downstream enabling role;
      identifies hubs with wide schedule impact \\[4pt]
Depth to sink $\delta_i$
    & $\max_{j:(i,j)\in E}(\delta_j+1)$
    & Distance to project end; complements
      $\ell_i$ for bidirectional positioning \\
\hline
\end{tabular}
\end{table}

\subsection{Extended Network Features (Group~B)}
\label{subsec:net_ext_features}

In-degree and out-degree only count an activity's \emph{direct}
neighbours, the activities immediately before or after it in
the precedence network. Three further features widen the view
by looking at the full chain of dependencies, not just the
immediate ones.

\paragraph{Transitive successors and predecessors.}
The \textit{number of transitive successors}
$n_{\mathrm{succ},i}$ counts every activity that depends on~$i$,
whether directly or through a chain of intermediate activities:
\begin{equation}
    n_{\mathrm{succ},i} = |\{j \in V : \text{there exists a
    directed path from } i \text{ to } j\}|
    \label{eq:trans_succ}
\end{equation}
If activity~$i$ has many transitive successors, a delay in~$i$
can cascade through a large part of the project.
The \textit{number of transitive predecessors}
$n_{\mathrm{pred},i}$ is the mirror image: it counts every
activity that must finish (directly or indirectly) before~$i$
can start:
\begin{equation}
    n_{\mathrm{pred},i} = |\{j \in V : \text{there exists a
    directed path from } j \text{ to } i\}|
    \label{eq:trans_pred}
\end{equation}
An activity with many transitive predecessors is a
convergence point: it is exposed to delays from many upstream
sources. Both counts are obtained by simple forward and backward
reachability passes over the topological order.

For example, consider a chain $A \to B \to C \to D$. Activity~$A$
has three transitive successors ($B$, $C$, $D$) and zero
transitive predecessors; activity~$D$ has zero transitive
successors and three transitive predecessors ($A$, $B$, $C$).
In contrast, their in- and out-degrees are both~1 (except at
the chain endpoints), so the transitive counts capture
long-range influence that the degree features miss.

\paragraph{Betweenness centrality.}
Some activities sit on almost every path through the project
network, while others lie on side branches that the project
could bypass. The \textit{betweenness centrality} quantifies
this by measuring the fraction of all source-to-sink paths
that pass through activity~$i$:
\begin{equation}
    \mathrm{bc}_i = \frac{\eta_{0 \to i} \cdot
    \eta_{i \to n{+}1}}{\eta_{0 \to n{+}1}}
    \label{eq:betweenness}
\end{equation}
Here $\eta_{0 \to i}$ is the number of distinct directed
paths from the project source (node~0) to activity~$i$, and
$\eta_{i \to n{+}1}$ is the number of paths from~$i$ to the
project sink (node~$n{+}1$). The product in the numerator
counts the source-to-sink paths that \emph{must} pass
through~$i$: any such path consists of a source-to-$i$
segment concatenated with an $i$-to-sink segment, and there
are $\eta_{0 \to i} \cdot \eta_{i \to n{+}1}$ such
combinations. Dividing by $\eta_{0 \to n{+}1}$ (the total
number of source-to-sink paths) normalises the count to
$[0, 1]$.

A value of $\mathrm{bc}_i = 1$ means every source-to-sink path
passes through~$i$; it is a mandatory bottleneck that the
project cannot route around. A value near~$0$ means the activity
sits on a minor branch: even if it is delayed, most paths
through the project are unaffected. Both $\eta$ counts are
computed efficiently by two DP passes (forward from the source
and backward from the sink) over the topological order.

\subsection{Duration Features (Group~C)}
\label{subsec:duration_features}

Three features describe how long each activity takes, both in
raw terms and relative to the rest of the instance:
\begin{align}
    \mathrm{dur}_i        &= d_i^d
    \label{eq:dur} \\[4pt]
    \mathrm{rel\_dur}_i   &= \frac{d_i^d}{d_{\max}},
    \qquad d_{\max} = \max_{j \in V \setminus \{0, n{+}1\}} d_j^d
    \label{eq:rel_dur} \\[4pt]
    \mathrm{rank}_i       &= \frac{|\{j \in V : d_j^d \leq d_i^d\}|}
                               {|V \setminus \{0, n{+}1\}|}
    \label{eq:dur_rank}
\end{align}
Here $d_i^d$ is the base processing time from the instance file,
the relative duration normalises it against the longest activity,
and the rank maps each activity to its percentile in the
instance-wide duration distribution (both in $[0,1]$; the rank
is less sensitive to outlier durations than the ratio).
Activities with similar durations occupy similar portions of the
scheduling horizon, so duration is a natural similarity signal
for grouping.

\subsection{Scheduling Urgency Features (Group~D)}
\label{subsec:sched_features}
Three features capture the degree to which the scheduling of
activity $i$ is constrained by its feasible time window
$[ES_i, LS_i]$, computed via the standard forward and backward
passes \parencite{Kelley1963}, and the level of resource
contention in its scheduling neighbourhood.

The \textit{slack}:
\begin{equation}
    \mathrm{slack}_i = LS_i - ES_i
    \label{eq:slack}
\end{equation}
measures how much the activity can be delayed without affecting
the project makespan. An activity with $\mathrm{slack}_i = 0$
lies on the critical path \parencite{Kelley1963} and must be scheduled
at exactly its earliest feasible start time. Slack has been widely
used as a priority rule in schedule generation schemes
\parencite{Kolisch1999} and is among the most informative features for
scheduling decisions \parencite{Hartmann1998}.

The \textit{pressure index}:
\begin{equation}
    \mathrm{pressure}_i = \frac{d_i^d}{LS_i - ES_i + 1}
    \label{eq:pressure}
\end{equation}
gives the fraction of the scheduling window that the activity's
base duration $d_i^d$ fills (the stochastic realisation
$d_i^{\pi}$ is not used here since it depends on the MIP
solution). When pressure is close to~1, the activity has almost
no room to shift; when it is close to~0, many feasible start
positions exist. The denominator $LS_i - ES_i + 1$ counts the
number of discrete start positions and avoids division by zero
for critical-path activities where $ES_i = LS_i$. Pressure adds
something slack alone misses: two activities can have the same
slack yet very different durations, and hence very different
scheduling flexibility.

The \textit{number of concurrent activities}:
\begin{equation}
    n_{\mathrm{conc},i} = \bigl|\bigl\{\, j \in V \setminus \{i\}
    : [ES_j,\, ES_j + d_j^d) \cap [ES_i,\, ES_i + d_i^d) \neq
    \emptyset \,\bigr\}\bigr|
    \label{eq:n_concurrent}
\end{equation}
counts the number of other real activities whose early-start
execution window overlaps with that of activity~$i$. The tight window $[ES_i,\, ES_i + d_i^d)$ is used rather
than the full feasible range $[ES_i,\, LS_i + d_i^d)$; in dense
RCPSP networks the latter produces a nearly constant count
because almost all activities overlap. With the tight window,
$n_{\mathrm{conc},i}$ reflects how crowded the earliest-possible
schedule actually is around activity~$i$. Where slack and
pressure describe an activity's own freedom, the concurrent count
shows how many neighbours are competing for the same resources in
the same time band.

\subsection{Timing Features (Group~E)}
\label{subsec:timing_features}

Four features describe the temporal position of each activity
within the scheduling horizon $[0, T]$, where
$T = \max_{j \in V} LS_j$:
\begin{align}
    \mathrm{rel\_es}_i   &= \frac{ES_i}{T}
    \label{eq:rel_es} \\[4pt]
    \mathrm{rel\_ls}_i   &= \frac{LS_i}{T}
    \label{eq:rel_ls} \\[4pt]
    \mathrm{cp}_i        &= \mathbb{1}[\mathrm{slack}_i = 0]
    \label{eq:critical_path} \\[4pt]
    \mathrm{float}_i     &= \frac{\mathrm{slack}_i}{T}
    \label{eq:float_ratio}
\end{align}
The relative earliest and latest start normalise the time
window endpoints to $[0,1]$, making them comparable across
instances with different horizons. The critical path indicator
is a binary feature that flags activities with zero slack.
The float ratio normalises slack by the horizon length,
capturing how much relative flexibility the activity has.

Unlike the urgency features, which describe local schedule
tightness, these four features pin each activity to a position
in the overall horizon. Two activities can share the same slack
and pressure yet sit at opposite ends of the project, facing
entirely different resource contention.

\subsection{Stochastic Resource Features (Groups~F--G)}
\label{subsec:stoch_features}
One joint feature (Group~F) summarises overall resource
demand, while eight per-resource features (Group~G) describe
how the stochastic workload varies across the $|S|$ scenarios.
Summarising stochastic behaviour with distributional statistics
is standard practice in the SRCPSP literature
\parencite{Herroelen2005, Lamas2016}.

The \textit{resource pressure}:
\begin{equation}
    \mathrm{rp}_i = \sum_{k \in K}
    \frac{\mu_{ik}}{R_k}
    \label{eq:resource_pressure}
\end{equation}
sums the mean workload $\mu_{ik}$
(equation~\ref{eq:mean_workload}) across all resources, each
divided by the corresponding capacity~$R_k$ so that a demand of
5~units on a 10-unit resource and a demand of 10~units on a
20-unit resource contribute equally. The result is a single
number reflecting how much total capacity the activity uses,
providing a bird's-eye complement to the per-resource detail
below.

For each resource $k \in K$, eight features describe the
distribution of the stochastic workload $w_{ik}^{\pi}$ across
the $|S|$ scenarios:

The \textit{mean workload}:
\begin{equation}
    \mu_{ik} = \frac{1}{|S|} \sum_{\pi \in S} w_{ik}^{\pi}
    \label{eq:mean_workload}
\end{equation}
captures the expected resource demand. The \textit{coefficient
of variation}:
\begin{equation}
    cv_{ik} = \frac{\sigma_{ik}}{\mu_{ik}},
    \qquad
    \sigma_{ik} = \sqrt{\frac{1}{|S|-1}
    \sum_{\pi \in S}(w_{ik}^{\pi} - \mu_{ik})^2}
    \label{eq:cv}
\end{equation}
measures relative variability \parencite{brown1998coefficient}.
Together, these two features encode the first two moments of
the workload distribution.

Six additional statistics complete the distributional
characterisation.

The \textit{interquartile range}:
\begin{equation}
    \mathrm{IQR}_{ik} = Q_{75}(w_{ik}^{\pi}) - Q_{25}(w_{ik}^{\pi})
    \label{eq:iqr}
\end{equation}
measures the spread of the middle $50\%$ of the workload
distribution. Unlike the standard deviation, IQR is robust to
outlier scenarios: a single extreme workload realisation does
not inflate it \parencite{Tukey1977}. This makes it a useful
complement to~$cv_{ik}$ for distinguishing activities whose
variability is driven by the bulk of scenarios from those whose
variability comes from one or two extreme cases.

The \textit{tail ratio}:
\begin{equation}
    \mathrm{tail}_{ik} = \frac{p_{90,ik}}{\mu_{ik}}
    \label{eq:tail_ratio}
\end{equation}
where $p_{90,ik}$ is the 90th percentile of $w_{ik}^{\pi}$
across the $|S|$ scenarios, quantifies how much heavier the
upper tail is relative to the mean. A tail ratio close to~1
indicates a symmetric or light-tailed distribution; values
well above~1 signal that worst-case scenarios demand
substantially more resources than the average case. For
$|S| = 10$, the 90th percentile corresponds to the second-largest
observation, so this feature captures near-worst-case behaviour.

The \textit{maximum workload}:
\begin{equation}
    \max_{ik} = \max_{\pi \in S}\; w_{ik}^{\pi}
    \label{eq:max_workload}
\end{equation}
records the highest resource demand observed across all
scenarios. This is the demand that determines whether the
activity can violate a resource capacity constraint in the
worst case.

The \textit{minimum workload}:
\begin{equation}
    \min_{ik} = \min_{\pi \in S}\; w_{ik}^{\pi}
    \label{eq:min_workload}
\end{equation}
records the lowest resource demand across all scenarios,
representing the best-case resource requirement.

The \textit{range}:
\begin{equation}
    \mathrm{range}_{ik} = \max_{ik} - \min_{ik}
    \label{eq:range}
\end{equation}
captures the total spread of the workload distribution from
best to worst case. Activities with a large range are highly
sensitive to the scenario realisation, while a range of zero
indicates deterministic demand on resource~$k$.

The \textit{nonzero probability}:
\begin{equation}
    P(\text{nz})_{ik} =
    \frac{|\{\pi \in S : w_{ik}^{\pi} > 0\}|}{|S|}
    \label{eq:pnz}
\end{equation}
gives the fraction of scenarios in which activity~$i$ requires
any amount of resource~$k$. An activity with
$P(\text{nz})_{ik} = 1$ always uses resource~$k$; one with a
low value uses it only in rare scenarios. This feature
distinguishes activities that consistently compete for a
resource from those that do so only occasionally, a distinction
that is invisible to the mean or variance alone.

For each resource~$k$, these eight features ($\mu_{ik}$,
$cv_{ik}$, $\mathrm{IQR}_{ik}$, $\mathrm{tail}_{ik}$,
$\max_{ik}$, $\min_{ik}$, $\mathrm{range}_{ik}$,
$P(\text{nz})_{ik}$) yield $8|K|$ resource columns in total.

\subsection{Feature Normalisation and Pruning}
\label{subsec:feat_norm}
Since features are measured in different units (time periods for
slack, dimensionless ratios for pressure and $cv$, resource units
for workload), direct distance-based clustering would be dominated
by features with large numerical ranges. All features are
therefore z-score scaled before clustering \parencite{Jain1999}:

\begin{equation}
    \tilde{x}_{if} = \frac{x_{if} - \bar{\mu}_f}{\bar{\sigma}_f}
    \label{eq:zscore}
\end{equation}

where $\bar{\mu}_f$ and $\bar{\sigma}_f$ are the mean and standard
deviation of feature $f$ across all activities in the instance.
After scaling, every feature has zero mean and unit variance,
so no single feature can dominate the Euclidean distance just
because its raw numbers happen to be large.

Features with pairwise Pearson correlation $|r| \geq
\theta_{\mathrm{corr}}$ are additionally pruned to avoid
redundancy, following the correlation-based feature selection
approach of \parencite{Hall1999}. Among the eight per-resource
statistics, empirical correlation analysis on PSPLIB instances
with $|K|=4$ and $|S|=10$ scenarios revealed several highly
correlated pairs: IQR with mean workload
($|r| = 0.95$--$1.00$), tail ratio with $cv$
($|r| = 0.96$--$0.97$), and $P(\text{nonzero})$ with mean
($|r| = 0.86$). The IQR--mean correlation is consistent with
the known relationship for beta-distributed variables
\parencite{Lamas2016}: when workloads are scaled by the same
factor, the IQR scales proportionally. The tail ratio
collapses to $cv$ because with $|S|=10$ scenarios the 90th
percentile approximates the maximum.

All $18 + 8|K|$ candidate features are nevertheless computed
and kept in the feature matrix. A runtime pruning step
($\theta_{\mathrm{corr}} = 0.95$) drops highly correlated
columns on a per-instance basis before clustering, so the
pipeline stays clean regardless of which features the user
selects. The comprehensive ablation
(Chapter~\ref{chap:experiments},
Section~\ref{sec:ablation}) confirms on all~50
features for $|K|=4$ that the pruning handles the redundancy
well.

The complete feature extraction procedure is given in
Algorithm~\ref{alg:features}.
% ── Algorithm 2 ──────────────────────────────────────────────
\begin{algorithm}[ht]
\caption{Comprehensive feature extraction for activity grouping}
\label{alg:features}
\begin{algorithmic}[1]
\State Compute precedence levels $\ell_i$ via forward
       longest-path DP (Algorithm~\ref{alg:prec_level_dp})
\State Compute depths $\delta_i$ via backward longest-path DP
       (Algorithm~\ref{alg:backward_dp})
\State Compute transitive reachability counts
       $n_{\mathrm{succ},i}$, $n_{\mathrm{pred},i}$
       and betweenness centrality $\mathrm{bc}_i$
\State Compute early-start windows
       $[ES_i,\, ES_i + d_i^d)$ for all $i$
\For{$i = 1$ \textbf{to} $n$}
    \State \textbf{[A]} Network basic:
           $\ell_i$, $\mathrm{in\_deg}_i$,
           $\mathrm{out\_deg}_i$, $\delta_i$
    \State \textbf{[B]} Network extended:
           $n_{\mathrm{succ},i}$,
           $n_{\mathrm{pred},i}$, $\mathrm{bc}_i$
    \State \textbf{[C]} Duration:
           $d_i^d$, $d_i^d / d_{\max}$,
           $\mathrm{rank}(d_i^d)$
    \State \textbf{[D]} Urgency:
           $\mathrm{slack}_i$,
           $\mathrm{pressure}_i$,
           $n_{\mathrm{conc},i}$
    \State \textbf{[E]} Timing:
           $ES_i/T$, $LS_i/T$,
           $\mathbb{1}[\mathrm{slack}_i{=}0]$,
           $\mathrm{slack}_i/T$
    \State \textbf{[F]} Resource joint:
           $\mathrm{rp}_i \leftarrow
           \sum_{k} \mu_{ik}/R_k$
    \For{$k = 1$ \textbf{to} $|K|$}
        \State \textbf{[G]} Per-resource stats:
               $\mu_{ik}$, $cv_{ik}$,
               $\mathrm{IQR}_{ik}$,
               $p_{90,ik}/\mu_{ik}$,
               $\max_{ik}$, $\min_{ik}$,
               $\mathrm{range}_{ik}$,
               $P(\text{nz})_{ik}$
    \EndFor
    \State Assemble $\mathbf{x}_i \in \mathbb{R}^{18+8|K|}$
           from groups A--G
\EndFor
\State Drop all-zero and constant columns
\State Z-score scale:
       $\tilde{x}_{if} = (x_{if}-\bar{\mu}_f)/\bar{\sigma}_f$
\State \Return Scaled feature matrix $\tilde{X}$
\end{algorithmic}
\end{algorithm}

\subsection{Default Feature Subset}
\label{subsec:default_features}
Although the full candidate pool contains $18 + 8|K|$ features
($50$ for $|K|=4$), a systematic ablation on PSPLIB j60
instances with $|K|=4$ (Chapter~\ref{chap:experiments},
Section~\ref{sec:ablation}) shows that using only the four
network position features (Group~A) achieves the best
trade-off between feasibility and solution quality. Across 32
configurations spanning 14 feature sub-groups and three
experimental phases, no configuration achieves both a lower gap
and equal or higher feasibility than the four network features. This result is consistent with the curse of
dimensionality in Ward clustering \parencite{Bellman1954}: adding
features from Groups~B--G increases the dimensionality from 4 to
up to 50, diluting the contribution of the structurally most
informative features (precedence level and depth) with noisier
signals.

Based on this evidence, the default feature subset restricts
clustering to the four network features:
\begin{equation}
    \mathbf{x}_i^{\mathrm{default}} =
    [\ell_i,\;\mathrm{in\_deg}_i,\;
     \mathrm{out\_deg}_i,\;\delta_i]
    \label{eq:default_features}
\end{equation}
This default can be overridden via the \texttt{feature\_cols}
parameter, allowing experiments with arbitrary feature subsets
from the full 50-feature candidate pool.

\subsubsection*{Summary of All Candidate Features}

Table~\ref{tab:all_features} consolidates the complete
candidate feature pool for reference.

\begin{table}[htbp]
\centering
\renewcommand{\arraystretch}{1.30}
\caption{Complete candidate feature pool ($18 + 8|K|$ features).
         Groups marked with $^{\dagger}$ contribute $|K|$
         features each. The default subset (marked $\star$)
         consists of Group~A only.}
\label{tab:all_features}
\begin{tabular}{c l r l l}
\hline
\textbf{Grp} & \textbf{Category} & $p$ &
\textbf{Features} & \textbf{Eq.} \\
\hline
A$^{\star}$ & Network basic
    & 4 & $\ell_i$, $\mathrm{in\_deg}_i$,
          $\mathrm{out\_deg}_i$, $\delta_i$
    & \eqref{eq:prec_level}--\eqref{eq:longest_path_sink} \\
B & Network extended
    & 3 & $n_{\mathrm{succ}}$, $n_{\mathrm{pred}}$,
          $\mathrm{bc}_i$
    & \eqref{eq:trans_succ}--\eqref{eq:betweenness} \\
C & Duration
    & 3 & $d_i^d$, $d_i^d/d_{\max}$,
          $\mathrm{rank}_i$
    & \eqref{eq:dur}--\eqref{eq:dur_rank} \\
D & Urgency
    & 3 & slack, pressure, $n_{\mathrm{conc}}$
    & \eqref{eq:slack}--\eqref{eq:n_concurrent} \\
E & Timing
    & 4 & rel.\ ES/LS, $\mathrm{cp}_i$,
          float ratio
    & \eqref{eq:rel_es}--\eqref{eq:float_ratio} \\
F & Resource joint
    & 1 & $\mathrm{rp}_i$
    & \eqref{eq:resource_pressure} \\
G$^{\dagger}$ & Per-resource stats
    & $8|K|$ & $\mu_{ik}$, $cv_{ik}$,
               IQR, tail, max, min,
               range, $P(\text{nz})$
    & \eqref{eq:mean_workload}--\eqref{eq:pnz} \\
\hline
 & \textbf{Total} & $\mathbf{18{+}8|K|}$ & & \\
\hline
\end{tabular}
\end{table}

% ------------------------------------------------------------
\section{Activity Grouping}
\label{sec:grouping}
% ------------------------------------------------------------
Given the scaled feature matrix $\tilde{X} \in \mathbb{R}^{n
\times p}$ (where $p$ depends on the selected feature subset,
defaulting to $p=4$ for the network-only configuration), the
grouping pipeline partitions the real activities into an ordered
list of groups $\mathcal{G} = [G_1, G_2, \ldots, G_m]$ satisfying
two properties: (i)~no group contains two activities connected by
a precedence arc, and (ii)~groups are ordered topologically so
that $G_g$ is always processed before $G_{g+1}$ in the R\&F loop.
\subsection{Precedence-Safe Partitioning}
\label{subsec:prec_safety}
A fundamental constraint on any grouping is that activities fixed
simultaneously as binary variables in one MIP subproblem must have
no precedence relationship between them. If activity $j$ is a
successor of activity $i$ and both are treated as binary in the
same subproblem, Gurobi may schedule $j$ before $i$, violating
the precedence constraint~(\ref{eq:precedence}) and potentially
rendering the subproblem infeasible or producing an invalid
schedule.

This constraint is guaranteed by restricting clustering to within
precedence levels. Activities $i$ and $j$ satisfy $\ell_i =
\ell_j$ if and only if neither is an ancestor of the other in
the precedence network, which is a direct consequence of the longest-path
DP in equation~(\ref{eq:prec_level}). The activities are therefore
first partitioned into levels $L_q = \{i \in V : \ell_i = q\}$
for $q = 1, \ldots, q_{\max}$, and Ward clustering is applied
independently within each level.

\subsection{Ward Hierarchical Clustering}
\label{subsec:ward}
Within each precedence level $L_q$, Ward hierarchical clustering
\parencite{Ward1963} is applied to the feature submatrix
$\tilde{X}[L_q]$. Ward's method is an agglomerative algorithm
that starts with each activity in its own cluster and
iteratively merges the pair whose combination produces the
smallest increase in total within-cluster variance
\parencite{Murtagh2012}:

\begin{equation}
    \Delta(A, B) = \frac{|A| \cdot |B|}{|A| + |B|}
    \cdot \|\boldsymbol{\mu}_A - \boldsymbol{\mu}_B\|_2^2
    \label{eq:ward}
\end{equation}

where $\boldsymbol{\mu}_A$ and $\boldsymbol{\mu}_B$ are the
centroids of clusters $A$ and $B$ respectively. The weighting
factor $|A||B|/(|A|+|B|)$ penalises merging large clusters more
heavily than small ones, which prevents the formation of
unbalanced clusters where one large group dominates the variance.
This produces a dendrogram, as in Figure~\ref{fig:dendrogram_level1}, encoding all possible groupings from
$|L_q|$ singletons down to a single cluster.

Ward's method suits this setting better than the obvious
alternatives. K-Means \parencite{MacQueen1967} would require
the cluster count~$k$ upfront, whereas the dendrogram encodes
every possible~$k$ at once and the optimal~$k^*$ can be read
off via the silhouette criterion below. Ward clustering is also
deterministic: the same feature matrix always yields the same
dendrogram, which makes results reproducible across
experimental runs. Finally, among the standard linkage criteria
it tends to produce the most compact, well-separated clusters
\parencite{Murtagh2012}, exactly the property needed when the
goal is to group activities with similar scheduling profiles.
\input{figures/dendogram}

\subsection{Optimal Number of Groups via Silhouette Score}
\label{subsec:silhouette}
The dendrogram is cut at the number of clusters $k^*$ that maximises
the average silhouette score \parencite{rousseeuw1987}:
\begin{equation}
    \bar{s}(k) = \frac{1}{|L_q|} \sum_{i \in L_q} s_i,
    \qquad
    s_i = \frac{b_i - a_i}{\max(a_i,\, b_i)}
    \label{eq:silhouette}
\end{equation}
where $a_i$ is the mean distance from activity $i$ to all other
activities in its cluster (cohesion) and $b_i$ is the mean distance
to all activities in the nearest other cluster (separation). The
silhouette score $s_i \in [-1, +1]$, with values near $+1$
indicating a well-placed activity and values near $-1$ indicating
a misclassified one. The optimal cut is:
\begin{equation}
    k^* = \arg\max_{k \in \{2, \ldots, |L_q|-1\}} \bar{s}(k)
    \label{eq:kstar}
\end{equation}
If the best silhouette score is below $0.10$, indicating that no
meaningful cluster structure exists, all activities in the level
are kept as one group.

\subsection{Merging Adjacent Groups}
\label{subsec:merge}
After Ward clustering assigns activities at each precedence
level to $k^*$ clusters, the grouping pipeline applies a
post-processing step that merges consecutive groups whose
feature centroids are sufficiently close. Since the silhouette
criterion picks $k^*$ independently within each level, two
clusters at adjacent levels can end up in separate groups even
though their activities are nearly identical in feature space.
Solving separate subproblems for such groups wastes time
without adding diversification to the search.

Consider two adjacent groups $G_g$ and $G_{g+1}$ in the
ordered sequence $\mathcal{G} = [G_1, \ldots, G_m]$. Their
feature centroids are defined as:

\begin{equation}
    \boldsymbol{\mu}_{G_g}
    = \frac{1}{|G_g|} \sum_{i \in G_g} \tilde{\mathbf{x}}_i,
    \qquad
    \boldsymbol{\mu}_{G_{g+1}}
    = \frac{1}{|G_{g+1}|} \sum_{i \in G_{g+1}}
      \tilde{\mathbf{x}}_i
    \label{eq:centroids}
\end{equation}

where $\tilde{\mathbf{x}}_i \in \mathbb{R}^p$ is the z-score
scaled, correlation-pruned feature vector of activity $i$.
When the centroid distance $\|\boldsymbol{\mu}_{G_g} -
\boldsymbol{\mu}_{G_{g+1}}\|_2$ is small, the two groups'
activities look alike in feature space, and fixing them
together in one subproblem gives the solver a richer
combinatorial context than two separate solves would (cf.\
Section~\ref{sec:grouprf_overview}).

\subsubsection{Merge Condition}
\label{subsubsec:merge_condition}

A merge of $G_g$ and $G_{g+1}$ is performed if and only if
both of the following conditions hold simultaneously.

\paragraph{Condition 1: Centroid proximity.}
\begin{equation}
    \|\boldsymbol{\mu}_{G_g} - \boldsymbol{\mu}_{G_{g+1}}\|_2
    < \theta_{\mathrm{merge}}
    \label{eq:merge_cond}
\end{equation}

where $\theta_{\mathrm{merge}} = 2.5$ is the merge threshold
in the normalised feature space. Since all features are
z-score scaled to zero mean and unit variance
(equation~\ref{eq:zscore}), this threshold is dimensionless
and scale-invariant across instances. A value of $2.5$
corresponds roughly to a moderate centroid separation:
groups whose centroids differ by more than $2.5$ standard
deviations in Euclidean distance are considered structurally
distinct and are kept separate.

\paragraph{Condition 2: Precedence safety.}
No activity in $G_{g+1}$ may be a direct or transitive
predecessor of any activity in $G_g$. Formally:

\begin{equation}
    \nexists\; i \in G_{g+1},\; j \in G_g
    \;\text{ such that }\;
    j \text{ is reachable from } i \text{ in } (V, E)
    \label{eq:merge_safety}
\end{equation}
This condition is necessary because merging $G_g$ and
$G_{g+1}$ means their activities will all be fixed
simultaneously as binary variables in one MIP subproblem, as
explained in Section~\ref{subsec:prec_safety}.
Since the groups are already ordered topologically by
construction, the safety condition reduces to checking
whether any backward arc exists in the transitive closure
$E^+$. This closure is computed once before the merge loop
and reused for all subsequent checks
\parencite{Demeulemeester2002}.

\subsubsection{Merge Procedure}
\label{subsubsec:merge_procedure}

The merge is applied in a repeated forward pass over the
group list until no further merges are possible, as
described in Algorithm~\ref{alg:merge}. The outer
\textbf{while} loop is necessary because merging $G_g$
with $G_{g+1}$ updates the centroid of the resulting
combined group, which may then satisfy the merge condition
with the next adjacent group $G_{g+1}$ (formerly $G_{g+2}$).
A single forward pass is therefore not guaranteed to find
all possible merges; the loop terminates only when a
complete pass produces no new merges.

\begin{algorithm}[ht]
\caption{Adjacent group merge post-processing}
\label{alg:merge}
\begin{algorithmic}[1]
\State Compute transitive closure $E^+$ of $(V, E)$
\State $\mathrm{changed} \leftarrow \mathrm{true}$
\While{$\mathrm{changed}$}
    \State $\mathrm{changed} \leftarrow \mathrm{false}$
    \For{$g = 1$ \textbf{to} $|\mathcal{G}| - 1$}
        \State $d \leftarrow
               \|\boldsymbol{\mu}_{G_g} -
                 \boldsymbol{\mu}_{G_{g+1}}\|_2$
        \State $\mathrm{safe} \leftarrow
               \nexists\; i \!\in\! G_{g+1},\;
               j \!\in\! G_g : (i,j) \in E^+$
        \If{$d < \theta_{\mathrm{merge}}$
            \textbf{and} $\mathrm{safe}$}
            \State $G_g \leftarrow G_g \cup G_{g+1}$,\;
                   recompute $\boldsymbol{\mu}_{G_g}$
            \State Remove $G_{g+1}$ from $\mathcal{G}$
            \State $\mathrm{changed} \leftarrow \mathrm{true}$
        \EndIf
    \EndFor
\EndWhile
\State \Return $\mathcal{G} = [G_1, \ldots, G_m]$
\end{algorithmic}
\end{algorithm}
\subsubsection{Worked Example}
\label{subsubsec:merge_example}

To make the procedure concrete, consider a small instance
where level-by-level clustering produces four initial groups:

\begin{table}[ht]
\centering
\renewcommand{\arraystretch}{1.3}
\caption{Initial groups before merge post-processing}
\label{tab:merge_example}
\begin{tabular}{clcc}
\hline
\textbf{Group} &
\textbf{Activities} &
$\boldsymbol{\mu}$ (first two dims) &
\textbf{Level} \\
\hline
$G_1$ & $\{1, 2\}$ & $(0.1,\; 0.3)$ & 1 \\
$G_2$ & $\{3\}$    & $(0.2,\; 0.4)$ & 1 \\
$G_3$ & $\{4, 5\}$ & $(1.8,\; 2.1)$ & 2 \\
$G_4$ & $\{6\}$    & $(0.3,\; 0.5)$ & 2 \\
\hline
\end{tabular}
\end{table}

\noindent
\textbf{Pass 1, $g=1$:}
$d(G_1, G_2) = \|(0.1,0.3)-(0.2,0.4)\|_2 \approx 0.14$.
Both groups are at level 1, so no precedence violation
exists between them.
Condition satisfied: $G_1 \leftarrow \{1,2,3\}$ with
updated centroid $(0.13, 0.33)$.
List becomes $[G_1, G_3, G_4]$.

\noindent
\textbf{Pass 1, $g=2$:}
$d(G_3, G_4) = \|(1.8,2.1)-(0.3,0.5)\|_2 \approx 2.12 < 2.5$.
No activity in $G_4$ is a predecessor of any activity in
$G_3$.
Condition satisfied: $G_3 \leftarrow \{4,5,6\}$.
List becomes $[G_1, G_3]$.
\noindent
\textbf{Pass 2, $g=1$:}
$d(G_1, G_3) = \|(0.13,0.33)-(1.30,1.57)\|_2 \approx 1.73 < 2.5$.
The centroid proximity condition is satisfied. For the
safety condition, the check is whether any activity
$i \in G_3$ is a predecessor of any activity $j \in G_1$
in the transitive closure~$E^+$. Since activities at
level~2 have activities at level~1 as predecessors (not
the reverse), the edges in~$E^+$ are of the form
$(j,i)$ with $j \in G_1$, $i \in G_3$, which does not
violate equation~(\ref{eq:merge_safety}). The safety
condition passes, and the merge proceeds, leaving
$\mathcal{G} = [\{1,2,3,4,5,6\}]$ as a single group.

In practice with $n = 30$, the centroid distance between
groups at different precedence levels is typically large
($d > 2.5$), preventing over-merging across levels.
The complete grouping procedure is summarised in
Algorithm~\ref{alg:grouping}.
% ------------------------------------------------------------
\section{Group-Based Relax-and-Fix}
\label{sec:group_rf_decomp}
% ------------------------------------------------------------
Given the ordered group list $\mathcal{G} = [G_1, \ldots, G_m]$,
the group-based R\&F proceeds analogously to the baseline but
processes one group per iteration instead of one activity.
\subsection{Three Variable States}
\label{subsec:three_states}
In each subproblem for group $G_g$, the decision variables are
assigned one of three states:
\begin{itemize}
    \item \textbf{Fixed:} Activities in previously processed groups
          $G_1, \ldots, G_{g-1}$ whose start times have been
          determined. Their binary variables are pinned:
          $x^B_{i, S_i} = 1$, $x^B_{it} = 0$ for $t \neq S_i$,
          $x^C_{it} = 0$.
    \item \textbf{Binary:} Activities in the current group $G_g$,
          treated as binary integer variables within their time
          windows: $x^B_{it} \in \{0,1\}$ for $t \in \widetilde{W}_i$.
    \item \textbf{Relaxed:} Activities in future groups
          $G_{g+1}, \ldots, G_m$ whose start times have not yet
          been determined. Their variables are relaxed to continuous:
          $x^C_{jt} \in [0,1]$.
\end{itemize}
By setting $x^B_{it} = 0$ for all fixed and relaxed activities
and $x^C_{it} = 0$ for all binary and fixed activities, the
effective start variable $x^{\mathrm{eff}}_{it} = x^B_{it} + x^C_{it}$
correctly reflects the state of each activity in every subproblem.
\subsection{Slack-Based Window Sizing}
\label{subsec:slack_window}
The window sizing strategy also differs from the baseline.
The baseline draws a random half-width
$w \sim \mathrm{Uniform}\{1, 10\}$ independently for each
activity subproblem, which can produce unnecessarily tight windows
(causing infeasible subproblems) or unnecessarily wide ones
(providing no improvement over the SSGS warm-start). The
group-based approach replaces this with a deterministic,
slack-proportional window width:

\begin{equation}
    w_i = \max\!\left(2,\;
    \left\lfloor \alpha_w \cdot (LS_i - ES_i) \right\rfloor\right),
    \quad \alpha_w = 0.15
    \label{eq:slack_window}
\end{equation}

The resulting time window for activity $i$ is centred on its
current schedule position $S'_i$:

\begin{equation}
    \widetilde{W}_i = \left[
    \max\{ES_i,\; S'_i - w_i\},\;
    \min\{LS_i,\; S'_i + w_i\}
    \right]
    \label{eq:group_window}
\end{equation}

Activities with large slack receive wider windows, giving the
solver more freedom to improve their placement. Critical-path
activities ($\mathrm{slack}_i = 0$) receive the minimum window
of~2, ensuring a minimal search neighbourhood around the
warm-start position.

\paragraph{Choice of $\alpha_w = 0.15$.}
The window fraction $\alpha_w$ defines a local search neighbourhood
around the warm-start position of each activity: the solver may
shift activity~$i$ by at most $w_i = \lfloor \alpha_w \cdot
\mathrm{slack}_i \rfloor$ periods from its current placement
$S'_i$. This is structurally analogous to the
\emph{destruction rate} in Large Neighbourhood Search~(LNS),
which controls what fraction of a solution is freed for
re-optimisation in each iteration \parencite{Pisinger2010}.
For combinatorial scheduling problems, destruction rates in the
range $[0.10, 0.30]$ consistently yield the best trade-off
between subproblem tractability and improvement potential
\parencite{Ropke2006, Pisinger2010}. Rates below $0.10$ produce
neighbourhoods too small to contain improving solutions; rates
above $0.30$ generate subproblems whose complexity approaches
that of the full model, negating the benefit of decomposition.
The value $\alpha_w = 0.15$ places the search radius conservatively
within this range, which is appropriate given that the group-based
decomposition already exposes $|G_g|$ activities simultaneously
--- a richer combinatorial structure per subproblem than
single-activity LNS.

\paragraph{Two-stage fallback.}
If the narrow-window subproblem yields no feasible solution
(typically because the restricted domain excludes all feasible
start-time combinations for the group), the solver retries
with the full feasible range $\widetilde{W}_i = [ES_i,\, LS_i]$
for all $i \in G_g$. This mirrors the adaptive neighbourhood
expansion used in LNS frameworks \parencite{Ropke2006}: when a
small neighbourhood fails, the search radius is increased.
A more gradual geometric expansion (e.g.\
$\alpha_w \in \{0.15, 0.30, 0.60, 1.0\}$), as suggested by the
progressive widening strategies of \textcite{Helber2010}, was not
adopted because the single-pass, stop-on-first-feasible
termination strategy
(Section~\ref{subsec:stopping}) makes total wall-clock time
across all groups the binding constraint: each intermediate
widening step that fails to find a solution consumes up to
$T_{\mathrm{sub}}$ seconds without advancing the schedule,
whereas the full-window fallback maximises the probability of
resolving the group in at most two MIP solves.

\subsection{Termination Strategy}
\label{subsec:stopping}
Three termination conditions govern the R\&F loop, checked in
order of priority after each group subproblem.

\paragraph{Stop on first feasible.}
When \texttt{stop\_on\_first\_feasible} is enabled (default), the
loop terminates as soon as any group subproblem returns a feasible
integer solution, i.e.\ a solution in which all effective start
variables satisfy $x^{\mathrm{eff}}_{it} \in \{0, 1\}$ for every
activity~$i$. Since the group-based decomposition fixes several related
activities at once, the first feasible solution typically
captures most of the attainable improvement; processing
further groups yields diminishing returns relative to the
additional computation time \parencite{pochet2006}.

\paragraph{Early stopping on consecutive non-improvement.}
The parameter \texttt{max\_no\_improve} (default~$0$, i.e.\
disabled) provides an alternative termination criterion for
configurations where \texttt{stop\_on\_first\_feasible} is
disabled. Once the first integral solution has been found, a
counter tracks consecutive groups that fail to improve the best
makespan. If this counter reaches \texttt{max\_no\_improve}, the
loop terminates.

\paragraph{Infeasible group fallback.}
If a group subproblem is infeasible after both the narrow-window
and full-window solves, the activities in $G_g$ are fixed at
their start times in the current warm-start schedule~$S'$, and
the loop proceeds to the next group. Fixing at~$S'$ rather than
at the original SSGS solution is preferable because $S'$ already
incorporates all improvements accepted in earlier groups
$G_1, \ldots, G_{g-1}$ \parencite{bigler2024multisite}.

\paragraph{Time limit.}
The overall time limit $T_{\max}$ provides a final termination
condition: if the cumulative wall-clock time exceeds $T_{\max}$
before any group produces a feasible solution, the algorithm
returns the best schedule found so far (which may be the SSGS
warm-start itself).

\medskip
\noindent
The complete group-based procedure, combining feature extraction
(Algorithm~\ref{alg:features}), activity grouping
(Algorithm~\ref{alg:grouping}), and the R\&F loop
(Algorithm~\ref{alg:group_rf}), is summarised in the two
pseudocode listings that follow.
% ── Algorithm 3 ──────────────────────────────────────────────
\begin{algorithm}[ht]
\caption{ML-based activity grouping}
\label{alg:grouping}
\begin{algorithmic}[1]
\State Restrict $\tilde{X}$ to selected feature subset
       (default: network features only,
        equation~\ref{eq:default_features})
\State Remove features with
       $|\mathrm{corr}(f_a,f_b)|\geq\theta_{\mathrm{corr}}$
       from $\tilde{X}$
\State Partition $V$ into levels
       $L_q=\{i\in V:\ell_i=q\}$,
       $q=1,\ldots,q_{\max}$
\For{$q = 1$ \textbf{to} $q_{\max}$}
    \State Build Ward dendrogram
           $Z=\mathrm{linkage}(\tilde{X}[L_q],\;\mathrm{ward})$
    \State $k^*\leftarrow\arg\max_k\,\bar{s}(k)$
           where $\bar{s}(k)=\frac{1}{|L_q|}\sum_{i\in L_q}s_i$
    \State Cut $Z$ into $k^*$ clusters and append to
           $\mathcal{G}$
\EndFor
\State $\mathrm{changed}\leftarrow\mathrm{true}$
\While{$\mathrm{changed}$}
    \State $\mathrm{changed}\leftarrow\mathrm{false}$
    \For{$g=1$ \textbf{to} $|\mathcal{G}|-1$}
        \If{$\|\boldsymbol{\mu}_{G_g}-
              \boldsymbol{\mu}_{G_{g+1}}\|_2
              <\theta_{\mathrm{merge}}$
              and no precedence violation}
            \State $G_g\leftarrow G_g\cup G_{g+1}$,
                   remove $G_{g+1}$
            \State $\mathrm{changed}\leftarrow\mathrm{true}$
        \EndIf
    \EndFor
\EndWhile
\State \Return $\mathcal{G}=[G_1,\ldots,G_m]$
\end{algorithmic}
\end{algorithm}
% ── Algorithm 8 ──────────────────────────────────────────────
\begin{algorithm}[H]
\caption{Procedure of the proposed group-based
         activity-oriented R\&F}
\label{alg:group_rf}
\begin{algorithmic}[1]
\State Determine greedy solution $S$ via SSGS
       (Algorithm~\ref{alg:ssgs})
\State Compute $\tilde{X}$ using Algorithm~\ref{alg:features}
\State Compute $\mathcal{G}=[G_1,\ldots,G_m]$
       using Algorithm~\ref{alg:grouping}
\State $\mathrm{no\_improve} \leftarrow 0$;\;
       $\mathrm{found\_integral} \leftarrow \mathrm{false}$
\For{$g=1$ \textbf{to} $|\mathcal{G}|$}
    \For{$i\in G_g$}
        \State $w_i \leftarrow \max(2,\,\lfloor \alpha_w \cdot
               (LS_i-ES_i)\rfloor)$
        \State $\widetilde{W}_i \leftarrow
               [\max\{ES_i,\,S'_i-w_i\},\;
                \min\{LS_i,\,S'_i+w_i\}]$
    \EndFor
    \State Solve SAA-FSRCPSP subproblem with narrow windows
    \If{no feasible solution}
        \State Set $\widetilde{W}_i=[ES_i,\,LS_i]$
               for all $i\in G_g$; re-solve with full windows
    \EndIf
    \If{feasible solution $S'$ found}
        \State Fix $x^B_{i,S'_i}:=1$ for all $i \in G_g$;\;
               update $S \leftarrow S'$
        \If{$S'$ improves $S_{\mathrm{best}}$}
            \State $S_{\mathrm{best}} \leftarrow S'$;\;
                   $\mathrm{no\_improve} \leftarrow 0$
        \ElsIf{$\mathrm{found\_integral}$}
            \State $\mathrm{no\_improve} \leftarrow
                   \mathrm{no\_improve} + 1$
        \EndIf
        \State $\mathrm{found\_integral} \leftarrow \mathrm{true}$
        \If{\texttt{stop\_on\_first\_feasible}}
            \State \textbf{break}
        \EndIf
    \Else
        \State Fix $i\in G_g$ at current $S'_i$
               \Comment{fallback}
        \If{$\mathrm{found\_integral}$}
            $\mathrm{no\_improve} \leftarrow
            \mathrm{no\_improve} + 1$
        \EndIf
    \EndIf
    \If{$\mathrm{max\_no\_improve} > 0$ \textbf{and}
        $\mathrm{no\_improve} \geq \mathrm{max\_no\_improve}$}
        \State \textbf{break}
               \Comment{early stopping}
    \EndIf
\EndFor
\State \Return Best solution $S_{\mathrm{best}}$
\end{algorithmic}
\end{algorithm}
\noindent Algorithm~\ref{alg:group_rf} presents the core
group-based R\&F procedure. For scaling to larger instances
($n \geq 60$), two optional extensions, namely time aggregation and
per-group scenario reduction, are introduced in
Section~\ref{sec:scalability}, with the extended algorithm
given in Algorithm~\ref{alg:group_rf_extended}.
% ------------------------------------------------------------
\section{Scalability Extensions}
\label{sec:scalability}
% ------------------------------------------------------------
The group-based R\&F procedure described in
Section~\ref{sec:group_rf_decomp} produces strong results on
instances with $n \in \{10, 30\}$ activities, but encounters
computational limits on larger instances. For $n = 60$ with
$|K| = 4$ resources and $|S| = 10$ scenarios, the SAA-FSRCPSP
model contains approximately 557{,}000 rows, 280{,}000 columns,
and 61 million nonzeros. At this scale, Gurobi's presolve and
initial LP relaxation alone consume 60--120 seconds per
subproblem, leaving little time for branch-and-bound within a
practical per-subproblem time limit of 60 seconds.

Two independent techniques are introduced to shrink the
per-subproblem MIP without changing the R\&F loop itself:
time aggregation (Section~\ref{subsec:time_agg}), which
compresses the time axis, and per-group scenario reduction
(Section~\ref{subsec:scen_red}), which keeps only a subset
of stochastic scenarios in each subproblem. Both act as
preprocessing, transforming the instance before each solve
and undoing the transformation afterwards. A final verification
pass (Section~\ref{subsec:final_verify}) then checks that the
schedule is still feasible on the original, unreduced instance.
Figure~\ref{fig:scalability_pipeline} summarises this
three-phase flow.

% ============================================================
% TikZ diagram: Three-Phase Scalability Pipeline
% Reduce → Solve → Verify (compact horizontal chain)
%
% Include in Chapter 5 with:
%   % ============================================================
% TikZ diagram: Three-Phase Scalability Pipeline
% Reduce → Solve → Verify (compact horizontal chain)
%
% Include in Chapter 5 with:
%   % ============================================================
% TikZ diagram: Three-Phase Scalability Pipeline
% Reduce → Solve → Verify (compact horizontal chain)
%
% Include in Chapter 5 with:
%   \input{figures/scalability_pipeline.tex}
%
% Required in preamble:
%   \usepackage{tikz}
%   \usetikzlibrary{arrows.meta, positioning}
% ============================================================
\begin{figure}[ht]
\centering
\resizebox{0.92\textwidth}{!}{%
\begin{tikzpicture}[
    node distance=0.55cm,
    box/.style={rectangle, draw=black!40, fill=gray!8,
                minimum height=0.85cm, minimum width=2cm,
                align=center, font=\footnotesize,
                inner sep=4pt},
    arr/.style={-{Stealth[length=4pt]}, draw=black!55},
    lbl/.style={font=\scriptsize, text=black!50},
]

\node[box] (orig)
    {Instance\\$T,\;|S|,\;n$};

\node[box, draw=blue!50, fill=blue!6,
      right=of orig] (reduce)
    {Reduce\\{\scriptsize $T\!\to\!\lceil T/\alpha\rceil$,\;
      $|S|\!\to\!n_{\mathrm{keep}}$}};

\node[box, draw=teal!50, fill=teal!6,
      right=of reduce] (solve)
    {Solve\\{\scriptsize group R\&F}};

\node[box, draw=orange!50, fill=orange!6,
      right=of solve] (verify)
    {Verify\\{\scriptsize full LP, all $|S|$}};

\node[box, right=of verify] (out)
    {$S^*$ +\\certificate};

\draw[arr] (orig)   -- (reduce);
\draw[arr] (reduce) -- (solve);
\draw[arr] (solve)  -- (verify);
\draw[arr] (verify) -- (out);

\end{tikzpicture}
}% end resizebox
\caption{Three-phase scalability pipeline: reduce the instance
         dimensions, solve via group-based R\&F on the smaller
         model, verify feasibility on the full instance.}
\label{fig:scalability_pipeline}
\end{figure}

%
% Required in preamble:
%   \usepackage{tikz}
%   \usetikzlibrary{arrows.meta, positioning}
% ============================================================
\begin{figure}[ht]
\centering
\resizebox{0.92\textwidth}{!}{%
\begin{tikzpicture}[
    node distance=0.55cm,
    box/.style={rectangle, draw=black!40, fill=gray!8,
                minimum height=0.85cm, minimum width=2cm,
                align=center, font=\footnotesize,
                inner sep=4pt},
    arr/.style={-{Stealth[length=4pt]}, draw=black!55},
    lbl/.style={font=\scriptsize, text=black!50},
]

\node[box] (orig)
    {Instance\\$T,\;|S|,\;n$};

\node[box, draw=blue!50, fill=blue!6,
      right=of orig] (reduce)
    {Reduce\\{\scriptsize $T\!\to\!\lceil T/\alpha\rceil$,\;
      $|S|\!\to\!n_{\mathrm{keep}}$}};

\node[box, draw=teal!50, fill=teal!6,
      right=of reduce] (solve)
    {Solve\\{\scriptsize group R\&F}};

\node[box, draw=orange!50, fill=orange!6,
      right=of solve] (verify)
    {Verify\\{\scriptsize full LP, all $|S|$}};

\node[box, right=of verify] (out)
    {$S^*$ +\\certificate};

\draw[arr] (orig)   -- (reduce);
\draw[arr] (reduce) -- (solve);
\draw[arr] (solve)  -- (verify);
\draw[arr] (verify) -- (out);

\end{tikzpicture}
}% end resizebox
\caption{Three-phase scalability pipeline: reduce the instance
         dimensions, solve via group-based R\&F on the smaller
         model, verify feasibility on the full instance.}
\label{fig:scalability_pipeline}
\end{figure}

%
% Required in preamble:
%   \usepackage{tikz}
%   \usetikzlibrary{arrows.meta, positioning}
% ============================================================
\begin{figure}[ht]
\centering
\resizebox{0.92\textwidth}{!}{%
\begin{tikzpicture}[
    node distance=0.55cm,
    box/.style={rectangle, draw=black!40, fill=gray!8,
                minimum height=0.85cm, minimum width=2cm,
                align=center, font=\footnotesize,
                inner sep=4pt},
    arr/.style={-{Stealth[length=4pt]}, draw=black!55},
    lbl/.style={font=\scriptsize, text=black!50},
]

\node[box] (orig)
    {Instance\\$T,\;|S|,\;n$};

\node[box, draw=blue!50, fill=blue!6,
      right=of orig] (reduce)
    {Reduce\\{\scriptsize $T\!\to\!\lceil T/\alpha\rceil$,\;
      $|S|\!\to\!n_{\mathrm{keep}}$}};

\node[box, draw=teal!50, fill=teal!6,
      right=of reduce] (solve)
    {Solve\\{\scriptsize group R\&F}};

\node[box, draw=orange!50, fill=orange!6,
      right=of solve] (verify)
    {Verify\\{\scriptsize full LP, all $|S|$}};

\node[box, right=of verify] (out)
    {$S^*$ +\\certificate};

\draw[arr] (orig)   -- (reduce);
\draw[arr] (reduce) -- (solve);
\draw[arr] (solve)  -- (verify);
\draw[arr] (verify) -- (out);

\end{tikzpicture}
}% end resizebox
\caption{Three-phase scalability pipeline: reduce the instance
         dimensions, solve via group-based R\&F on the smaller
         model, verify feasibility on the full instance.}
\label{fig:scalability_pipeline}
\end{figure}


% ────────────────────────────────────────────────────────────
\subsection{Time Aggregation}
\label{subsec:time_agg}
% ────────────────────────────────────────────────────────────

Every variable in the SAA-FSRCPSP carries a time index:
$x^B_{it}$ and $x^C_{it}$ exist for each activity~$i$ and each
period~$t \in [ES_i, LS_i]$, while $r_{ikt\pi}$ exists for each
activity, resource, period, and scenario. Variable and constraint
counts therefore grow linearly with the horizon~$T$, and for
$n = 60$ instances with $T \approx 150$ the time axis is the
single largest contributor to model size.

Time aggregation replaces $T$ individual time periods with
$T_{\mathrm{agg}} = \lceil T / \alpha \rceil$ aggregated
\emph{buckets}, where $\alpha \in \mathbb{N}_{\geq 1}$ is the
aggregation factor. Each bucket $\tau$ represents $\alpha$
consecutive physical time periods:
\begin{equation}
    \mathrm{bucket}(\tau) =
    \{\alpha\tau,\; \alpha\tau + 1,\; \ldots,\;
      \alpha\tau + \alpha - 1\}
    \label{eq:bucket}
\end{equation}
Since every time-indexed tensor ($x^B$, $x^C$, $y$, $r$) shrinks
by a factor of $\alpha$ along the $t$-axis, a choice of
$\alpha = 5$ reduces the time dimension from 150 to 30, an $80\%$
reduction in the number of time-indexed variables and constraints.
The aggregated instance is structurally identical to the original,
with the following transformations applied to the instance data.
All quantities not listed below (in particular the workloads
$w_{ik}^{\pi}$, the precedence graph $(V, E)$, and the number of
activities $n$) remain unchanged, since total work is invariant
under time bucketing and the precedence structure does not depend
on time granularity.

Figure~\ref{fig:time_agg} illustrates the transformation for two
activities on a horizon of $T = 15$ periods with $\alpha = 5$.
The physical timeline is compressed from 15~periods to
3~buckets; the floor/ceiling rules for $ES$, $LS$, and $d$
used in the figure are formalised below.

% ============================================================
% TikZ diagram: Time Aggregation Example
% Shows physical timeline (T=15) compressed into buckets (alpha=5)
%
% Include in Chapter 5 with:
%   % ============================================================
% TikZ diagram: Time Aggregation Example
% Shows physical timeline (T=15) compressed into buckets (alpha=5)
%
% Include in Chapter 5 with:
%   % ============================================================
% TikZ diagram: Time Aggregation Example
% Shows physical timeline (T=15) compressed into buckets (alpha=5)
%
% Include in Chapter 5 with:
%   \input{figures/time_aggregation.tex}
%
% Required in preamble:
%   \usepackage{tikz}
%   \usetikzlibrary{arrows.meta, positioning, calc}
% ============================================================
\begin{figure}[ht]
\centering
\begin{tikzpicture}[
    >=Stealth,
    lbl/.style={font=\scriptsize},
    tlbl/.style={font=\tiny, text=black!60},
]

% ================================================================
%  (a) PHYSICAL TIMELINE — T = 15
% ================================================================
\def\W{12}            % total diagram width in cm
\def\dx{0.8}          % width per period = 12/15
\def\ya{0}            % y-origin for part (a)

\node[font=\small\bfseries, anchor=south west]
    at (0, \ya+2.6) {(a)\; Physical timeline \;($T = 15$)};

% ── Bucket background shading ────────────────────────────
\fill[blue!6]   (0,       \ya-1.4) rectangle (5*\dx, \ya+2.2);
\fill[blue!12]  (5*\dx,   \ya-1.4) rectangle (10*\dx, \ya+2.2);
\fill[blue!6]   (10*\dx,  \ya-1.4) rectangle (15*\dx, \ya+2.2);

% Bucket labels (top)
\node[font=\scriptsize, blue!60] at (2.5*\dx, \ya+2.0)
    {bucket $\tau{=}0$};
\node[font=\scriptsize, blue!60] at (7.5*\dx, \ya+2.0)
    {bucket $\tau{=}1$};
\node[font=\scriptsize, blue!60] at (12.5*\dx, \ya+2.0)
    {bucket $\tau{=}2$};

% ── Time axis ────────────────────────────────────────────
\draw[->, thick, black!50] (-0.15, \ya) -- (\W+0.4, \ya)
    node[right, lbl, black!50] {$t$};

% Period ticks
\foreach \t in {0,1,...,15} {
    \draw[black!30] (\t*\dx, \ya+0.08) -- (\t*\dx, \ya-0.08);
    \node[tlbl, below=1pt] at (\t*\dx, \ya-0.08) {\t};
}

% Bucket boundary lines (heavier)
\foreach \b in {0, 5, 10, 15} {
    \draw[blue!30, thin, dashed]
        (\b*\dx, \ya-1.4) -- (\b*\dx, \ya+2.2);
}

% ── Activity i (teal): ES=2, LS=8, d=3 ──────────────────
%    Show the activity bar at its ES position
\fill[teal!50, rounded corners=2pt]
    (2*\dx, \ya+0.5) rectangle (5*\dx, \ya+1.0);
\node[font=\scriptsize, text=white, font=\bfseries]
    at (3.5*\dx, \ya+0.75) {$i$\; ($d_i{=}3$)};

% ES and LS markers
\draw[teal!70, thick] (2*\dx, \ya+0.35) -- (2*\dx, \ya+1.15);
\draw[teal!70, thick] (8*\dx, \ya+0.35) -- (8*\dx, \ya+1.15);
\node[lbl, teal!80, above=1pt] at (2*\dx, \ya+1.15) {$ES_i{=}2$};
\node[lbl, teal!80, above=1pt] at (8*\dx, \ya+1.15) {$LS_i{=}8$};

% Feasible window bracket
\draw[teal!40, <->, thick]
    (2*\dx, \ya+0.35) -- (8*\dx, \ya+0.35);

% ── Activity j (violet): ES=7, LS=12, d=4 ───────────────
\fill[violet!50, rounded corners=2pt]
    (7*\dx, \ya-1.1) rectangle (11*\dx, \ya-0.6);
\node[font=\scriptsize, text=white, font=\bfseries]
    at (9*\dx, \ya-0.85) {$j$\; ($d_j{=}4$)};

% ES and LS markers
\draw[violet!70, thick] (7*\dx, \ya-0.45) -- (7*\dx, \ya-1.25);
\draw[violet!70, thick] (12*\dx, \ya-0.45) -- (12*\dx, \ya-1.25);
\node[lbl, violet!80, below=1pt] at (7*\dx, \ya-1.25)
    {$ES_j{=}7$};
\node[lbl, violet!80, below=1pt] at (12*\dx, \ya-1.25)
    {$LS_j{=}12$};

% Feasible window bracket
\draw[violet!40, <->, thick]
    (7*\dx, \ya-0.45) -- (12*\dx, \ya-0.45);

% ================================================================
%  Arrow between (a) and (b)
% ================================================================
\draw[->, line width=1.2pt, black!40]
    (6, \ya-1.9) -- (6, \ya-2.7)
    node[midway, right=3pt, font=\scriptsize, black!50]
    {aggregate by $\alpha{=}5$};

% ================================================================
%  (b) AGGREGATED TIMELINE — T_agg = 3
% ================================================================
\def\yb{-5.5}         % y-origin for part (b)
\def\bx{4.0}          % width per bucket = 12/3

\node[font=\small\bfseries, anchor=south west]
    at (0, \yb+2.6) {(b)\; Aggregated timeline
    \;($T_{\mathrm{agg}} = 3$)};

% ── Bucket background shading ────────────────────────────
\fill[blue!6]   (0,       \yb-1.4) rectangle (1*\bx, \yb+2.2);
\fill[blue!12]  (1*\bx,   \yb-1.4) rectangle (2*\bx, \yb+2.2);
\fill[blue!6]   (2*\bx,   \yb-1.4) rectangle (3*\bx, \yb+2.2);

% Bucket labels (top)
\node[font=\scriptsize, blue!60] at (0.5*\bx, \yb+2.0) {$\tau{=}0$};
\node[font=\scriptsize, blue!60] at (1.5*\bx, \yb+2.0) {$\tau{=}1$};
\node[font=\scriptsize, blue!60] at (2.5*\bx, \yb+2.0) {$\tau{=}2$};

% ── Time axis ────────────────────────────────────────────
\draw[->, thick, black!50] (-0.15, \yb) -- (3*\bx+0.4, \yb)
    node[right, lbl, black!50] {$\tau$};

% Bucket ticks
\foreach \t in {0,1,2,3} {
    \draw[black!40] (\t*\bx, \yb+0.08) -- (\t*\bx, \yb-0.08);
    \node[lbl, below=1pt, black!60] at (\t*\bx, \yb-0.08) {\t};
}

% Bucket boundary lines
\foreach \b in {0,1,2,3} {
    \draw[blue!30, thin, dashed]
        (\b*\bx, \yb-1.4) -- (\b*\bx, \yb+2.2);
}

% ── Activity i aggregated ────────────────────────────────
%    ES_agg = floor(2/5) = 0,  LS_agg = ceil(8/5) = 2,  d_agg = ceil(3/5) = 1
\fill[teal!50, rounded corners=2pt]
    (0*\bx, \yb+0.5) rectangle (1*\bx, \yb+1.0);
\node[font=\scriptsize, text=white, font=\bfseries]
    at (0.5*\bx, \yb+0.75) {$i$\; ($d_i^{\mathrm{agg}}{=}1$)};

% ES and LS markers
\draw[teal!70, thick] (0*\bx, \yb+0.35) -- (0*\bx, \yb+1.15);
\draw[teal!70, thick] (2*\bx, \yb+0.35) -- (2*\bx, \yb+1.15);
\node[lbl, teal!80, above=1pt] at (0*\bx, \yb+1.15)
    {$\lfloor 2/5 \rfloor{=}0$};
\node[lbl, teal!80, above=1pt] at (2*\bx, \yb+1.15)
    {$\lceil 8/5 \rceil{=}2$};

% Feasible window bracket
\draw[teal!40, <->, thick]
    (0*\bx, \yb+0.35) -- (2*\bx, \yb+0.35);

% ── Activity j aggregated ────────────────────────────────
%    ES_agg = floor(7/5) = 1,  LS_agg = ceil(12/5) = 3,  d_agg = ceil(4/5) = 1
\fill[violet!50, rounded corners=2pt]
    (1*\bx, \yb-1.1) rectangle (2*\bx, \yb-0.6);
\node[font=\scriptsize, text=white, font=\bfseries]
    at (1.5*\bx, \yb-0.85) {$j$\; ($d_j^{\mathrm{agg}}{=}1$)};

% ES and LS markers
\draw[violet!70, thick] (1*\bx, \yb-0.45) -- (1*\bx, \yb-1.25);
\draw[violet!70, thick] (3*\bx, \yb-0.45) -- (3*\bx, \yb-1.25);
\node[lbl, violet!80, below=1pt] at (1*\bx, \yb-1.25)
    {$\lfloor 7/5 \rfloor{=}1$};
\node[lbl, violet!80, below=1pt] at (3*\bx, \yb-1.25)
    {$\lceil 12/5 \rceil{=}3$};

% Feasible window bracket
\draw[violet!40, <->, thick]
    (1*\bx, \yb-0.45) -- (3*\bx, \yb-0.45);

\end{tikzpicture}
\caption{Time aggregation with factor $\alpha = 5$.
         Activities $i$ and $j$ are shown on the physical
         timeline~(a) and the aggregated timeline~(b).
         Coloured bars show each activity at its earliest start;
         double-headed arrows mark the feasible window $[ES, LS]$.
         Floor division maps $ES$ into buckets; ceiling division
         maps $LS$ and durations to preserve feasibility.}
\label{fig:time_agg}
\end{figure}

%
% Required in preamble:
%   \usepackage{tikz}
%   \usetikzlibrary{arrows.meta, positioning, calc}
% ============================================================
\begin{figure}[ht]
\centering
\begin{tikzpicture}[
    >=Stealth,
    lbl/.style={font=\scriptsize},
    tlbl/.style={font=\tiny, text=black!60},
]

% ================================================================
%  (a) PHYSICAL TIMELINE — T = 15
% ================================================================
\def\W{12}            % total diagram width in cm
\def\dx{0.8}          % width per period = 12/15
\def\ya{0}            % y-origin for part (a)

\node[font=\small\bfseries, anchor=south west]
    at (0, \ya+2.6) {(a)\; Physical timeline \;($T = 15$)};

% ── Bucket background shading ────────────────────────────
\fill[blue!6]   (0,       \ya-1.4) rectangle (5*\dx, \ya+2.2);
\fill[blue!12]  (5*\dx,   \ya-1.4) rectangle (10*\dx, \ya+2.2);
\fill[blue!6]   (10*\dx,  \ya-1.4) rectangle (15*\dx, \ya+2.2);

% Bucket labels (top)
\node[font=\scriptsize, blue!60] at (2.5*\dx, \ya+2.0)
    {bucket $\tau{=}0$};
\node[font=\scriptsize, blue!60] at (7.5*\dx, \ya+2.0)
    {bucket $\tau{=}1$};
\node[font=\scriptsize, blue!60] at (12.5*\dx, \ya+2.0)
    {bucket $\tau{=}2$};

% ── Time axis ────────────────────────────────────────────
\draw[->, thick, black!50] (-0.15, \ya) -- (\W+0.4, \ya)
    node[right, lbl, black!50] {$t$};

% Period ticks
\foreach \t in {0,1,...,15} {
    \draw[black!30] (\t*\dx, \ya+0.08) -- (\t*\dx, \ya-0.08);
    \node[tlbl, below=1pt] at (\t*\dx, \ya-0.08) {\t};
}

% Bucket boundary lines (heavier)
\foreach \b in {0, 5, 10, 15} {
    \draw[blue!30, thin, dashed]
        (\b*\dx, \ya-1.4) -- (\b*\dx, \ya+2.2);
}

% ── Activity i (teal): ES=2, LS=8, d=3 ──────────────────
%    Show the activity bar at its ES position
\fill[teal!50, rounded corners=2pt]
    (2*\dx, \ya+0.5) rectangle (5*\dx, \ya+1.0);
\node[font=\scriptsize, text=white, font=\bfseries]
    at (3.5*\dx, \ya+0.75) {$i$\; ($d_i{=}3$)};

% ES and LS markers
\draw[teal!70, thick] (2*\dx, \ya+0.35) -- (2*\dx, \ya+1.15);
\draw[teal!70, thick] (8*\dx, \ya+0.35) -- (8*\dx, \ya+1.15);
\node[lbl, teal!80, above=1pt] at (2*\dx, \ya+1.15) {$ES_i{=}2$};
\node[lbl, teal!80, above=1pt] at (8*\dx, \ya+1.15) {$LS_i{=}8$};

% Feasible window bracket
\draw[teal!40, <->, thick]
    (2*\dx, \ya+0.35) -- (8*\dx, \ya+0.35);

% ── Activity j (violet): ES=7, LS=12, d=4 ───────────────
\fill[violet!50, rounded corners=2pt]
    (7*\dx, \ya-1.1) rectangle (11*\dx, \ya-0.6);
\node[font=\scriptsize, text=white, font=\bfseries]
    at (9*\dx, \ya-0.85) {$j$\; ($d_j{=}4$)};

% ES and LS markers
\draw[violet!70, thick] (7*\dx, \ya-0.45) -- (7*\dx, \ya-1.25);
\draw[violet!70, thick] (12*\dx, \ya-0.45) -- (12*\dx, \ya-1.25);
\node[lbl, violet!80, below=1pt] at (7*\dx, \ya-1.25)
    {$ES_j{=}7$};
\node[lbl, violet!80, below=1pt] at (12*\dx, \ya-1.25)
    {$LS_j{=}12$};

% Feasible window bracket
\draw[violet!40, <->, thick]
    (7*\dx, \ya-0.45) -- (12*\dx, \ya-0.45);

% ================================================================
%  Arrow between (a) and (b)
% ================================================================
\draw[->, line width=1.2pt, black!40]
    (6, \ya-1.9) -- (6, \ya-2.7)
    node[midway, right=3pt, font=\scriptsize, black!50]
    {aggregate by $\alpha{=}5$};

% ================================================================
%  (b) AGGREGATED TIMELINE — T_agg = 3
% ================================================================
\def\yb{-5.5}         % y-origin for part (b)
\def\bx{4.0}          % width per bucket = 12/3

\node[font=\small\bfseries, anchor=south west]
    at (0, \yb+2.6) {(b)\; Aggregated timeline
    \;($T_{\mathrm{agg}} = 3$)};

% ── Bucket background shading ────────────────────────────
\fill[blue!6]   (0,       \yb-1.4) rectangle (1*\bx, \yb+2.2);
\fill[blue!12]  (1*\bx,   \yb-1.4) rectangle (2*\bx, \yb+2.2);
\fill[blue!6]   (2*\bx,   \yb-1.4) rectangle (3*\bx, \yb+2.2);

% Bucket labels (top)
\node[font=\scriptsize, blue!60] at (0.5*\bx, \yb+2.0) {$\tau{=}0$};
\node[font=\scriptsize, blue!60] at (1.5*\bx, \yb+2.0) {$\tau{=}1$};
\node[font=\scriptsize, blue!60] at (2.5*\bx, \yb+2.0) {$\tau{=}2$};

% ── Time axis ────────────────────────────────────────────
\draw[->, thick, black!50] (-0.15, \yb) -- (3*\bx+0.4, \yb)
    node[right, lbl, black!50] {$\tau$};

% Bucket ticks
\foreach \t in {0,1,2,3} {
    \draw[black!40] (\t*\bx, \yb+0.08) -- (\t*\bx, \yb-0.08);
    \node[lbl, below=1pt, black!60] at (\t*\bx, \yb-0.08) {\t};
}

% Bucket boundary lines
\foreach \b in {0,1,2,3} {
    \draw[blue!30, thin, dashed]
        (\b*\bx, \yb-1.4) -- (\b*\bx, \yb+2.2);
}

% ── Activity i aggregated ────────────────────────────────
%    ES_agg = floor(2/5) = 0,  LS_agg = ceil(8/5) = 2,  d_agg = ceil(3/5) = 1
\fill[teal!50, rounded corners=2pt]
    (0*\bx, \yb+0.5) rectangle (1*\bx, \yb+1.0);
\node[font=\scriptsize, text=white, font=\bfseries]
    at (0.5*\bx, \yb+0.75) {$i$\; ($d_i^{\mathrm{agg}}{=}1$)};

% ES and LS markers
\draw[teal!70, thick] (0*\bx, \yb+0.35) -- (0*\bx, \yb+1.15);
\draw[teal!70, thick] (2*\bx, \yb+0.35) -- (2*\bx, \yb+1.15);
\node[lbl, teal!80, above=1pt] at (0*\bx, \yb+1.15)
    {$\lfloor 2/5 \rfloor{=}0$};
\node[lbl, teal!80, above=1pt] at (2*\bx, \yb+1.15)
    {$\lceil 8/5 \rceil{=}2$};

% Feasible window bracket
\draw[teal!40, <->, thick]
    (0*\bx, \yb+0.35) -- (2*\bx, \yb+0.35);

% ── Activity j aggregated ────────────────────────────────
%    ES_agg = floor(7/5) = 1,  LS_agg = ceil(12/5) = 3,  d_agg = ceil(4/5) = 1
\fill[violet!50, rounded corners=2pt]
    (1*\bx, \yb-1.1) rectangle (2*\bx, \yb-0.6);
\node[font=\scriptsize, text=white, font=\bfseries]
    at (1.5*\bx, \yb-0.85) {$j$\; ($d_j^{\mathrm{agg}}{=}1$)};

% ES and LS markers
\draw[violet!70, thick] (1*\bx, \yb-0.45) -- (1*\bx, \yb-1.25);
\draw[violet!70, thick] (3*\bx, \yb-0.45) -- (3*\bx, \yb-1.25);
\node[lbl, violet!80, below=1pt] at (1*\bx, \yb-1.25)
    {$\lfloor 7/5 \rfloor{=}1$};
\node[lbl, violet!80, below=1pt] at (3*\bx, \yb-1.25)
    {$\lceil 12/5 \rceil{=}3$};

% Feasible window bracket
\draw[violet!40, <->, thick]
    (1*\bx, \yb-0.45) -- (3*\bx, \yb-0.45);

\end{tikzpicture}
\caption{Time aggregation with factor $\alpha = 5$.
         Activities $i$ and $j$ are shown on the physical
         timeline~(a) and the aggregated timeline~(b).
         Coloured bars show each activity at its earliest start;
         double-headed arrows mark the feasible window $[ES, LS]$.
         Floor division maps $ES$ into buckets; ceiling division
         maps $LS$ and durations to preserve feasibility.}
\label{fig:time_agg}
\end{figure}

%
% Required in preamble:
%   \usepackage{tikz}
%   \usetikzlibrary{arrows.meta, positioning, calc}
% ============================================================
\begin{figure}[ht]
\centering
\begin{tikzpicture}[
    >=Stealth,
    lbl/.style={font=\scriptsize},
    tlbl/.style={font=\tiny, text=black!60},
]

% ================================================================
%  (a) PHYSICAL TIMELINE — T = 15
% ================================================================
\def\W{12}            % total diagram width in cm
\def\dx{0.8}          % width per period = 12/15
\def\ya{0}            % y-origin for part (a)

\node[font=\small\bfseries, anchor=south west]
    at (0, \ya+2.6) {(a)\; Physical timeline \;($T = 15$)};

% ── Bucket background shading ────────────────────────────
\fill[blue!6]   (0,       \ya-1.4) rectangle (5*\dx, \ya+2.2);
\fill[blue!12]  (5*\dx,   \ya-1.4) rectangle (10*\dx, \ya+2.2);
\fill[blue!6]   (10*\dx,  \ya-1.4) rectangle (15*\dx, \ya+2.2);

% Bucket labels (top)
\node[font=\scriptsize, blue!60] at (2.5*\dx, \ya+2.0)
    {bucket $\tau{=}0$};
\node[font=\scriptsize, blue!60] at (7.5*\dx, \ya+2.0)
    {bucket $\tau{=}1$};
\node[font=\scriptsize, blue!60] at (12.5*\dx, \ya+2.0)
    {bucket $\tau{=}2$};

% ── Time axis ────────────────────────────────────────────
\draw[->, thick, black!50] (-0.15, \ya) -- (\W+0.4, \ya)
    node[right, lbl, black!50] {$t$};

% Period ticks
\foreach \t in {0,1,...,15} {
    \draw[black!30] (\t*\dx, \ya+0.08) -- (\t*\dx, \ya-0.08);
    \node[tlbl, below=1pt] at (\t*\dx, \ya-0.08) {\t};
}

% Bucket boundary lines (heavier)
\foreach \b in {0, 5, 10, 15} {
    \draw[blue!30, thin, dashed]
        (\b*\dx, \ya-1.4) -- (\b*\dx, \ya+2.2);
}

% ── Activity i (teal): ES=2, LS=8, d=3 ──────────────────
%    Show the activity bar at its ES position
\fill[teal!50, rounded corners=2pt]
    (2*\dx, \ya+0.5) rectangle (5*\dx, \ya+1.0);
\node[font=\scriptsize, text=white, font=\bfseries]
    at (3.5*\dx, \ya+0.75) {$i$\; ($d_i{=}3$)};

% ES and LS markers
\draw[teal!70, thick] (2*\dx, \ya+0.35) -- (2*\dx, \ya+1.15);
\draw[teal!70, thick] (8*\dx, \ya+0.35) -- (8*\dx, \ya+1.15);
\node[lbl, teal!80, above=1pt] at (2*\dx, \ya+1.15) {$ES_i{=}2$};
\node[lbl, teal!80, above=1pt] at (8*\dx, \ya+1.15) {$LS_i{=}8$};

% Feasible window bracket
\draw[teal!40, <->, thick]
    (2*\dx, \ya+0.35) -- (8*\dx, \ya+0.35);

% ── Activity j (violet): ES=7, LS=12, d=4 ───────────────
\fill[violet!50, rounded corners=2pt]
    (7*\dx, \ya-1.1) rectangle (11*\dx, \ya-0.6);
\node[font=\scriptsize, text=white, font=\bfseries]
    at (9*\dx, \ya-0.85) {$j$\; ($d_j{=}4$)};

% ES and LS markers
\draw[violet!70, thick] (7*\dx, \ya-0.45) -- (7*\dx, \ya-1.25);
\draw[violet!70, thick] (12*\dx, \ya-0.45) -- (12*\dx, \ya-1.25);
\node[lbl, violet!80, below=1pt] at (7*\dx, \ya-1.25)
    {$ES_j{=}7$};
\node[lbl, violet!80, below=1pt] at (12*\dx, \ya-1.25)
    {$LS_j{=}12$};

% Feasible window bracket
\draw[violet!40, <->, thick]
    (7*\dx, \ya-0.45) -- (12*\dx, \ya-0.45);

% ================================================================
%  Arrow between (a) and (b)
% ================================================================
\draw[->, line width=1.2pt, black!40]
    (6, \ya-1.9) -- (6, \ya-2.7)
    node[midway, right=3pt, font=\scriptsize, black!50]
    {aggregate by $\alpha{=}5$};

% ================================================================
%  (b) AGGREGATED TIMELINE — T_agg = 3
% ================================================================
\def\yb{-5.5}         % y-origin for part (b)
\def\bx{4.0}          % width per bucket = 12/3

\node[font=\small\bfseries, anchor=south west]
    at (0, \yb+2.6) {(b)\; Aggregated timeline
    \;($T_{\mathrm{agg}} = 3$)};

% ── Bucket background shading ────────────────────────────
\fill[blue!6]   (0,       \yb-1.4) rectangle (1*\bx, \yb+2.2);
\fill[blue!12]  (1*\bx,   \yb-1.4) rectangle (2*\bx, \yb+2.2);
\fill[blue!6]   (2*\bx,   \yb-1.4) rectangle (3*\bx, \yb+2.2);

% Bucket labels (top)
\node[font=\scriptsize, blue!60] at (0.5*\bx, \yb+2.0) {$\tau{=}0$};
\node[font=\scriptsize, blue!60] at (1.5*\bx, \yb+2.0) {$\tau{=}1$};
\node[font=\scriptsize, blue!60] at (2.5*\bx, \yb+2.0) {$\tau{=}2$};

% ── Time axis ────────────────────────────────────────────
\draw[->, thick, black!50] (-0.15, \yb) -- (3*\bx+0.4, \yb)
    node[right, lbl, black!50] {$\tau$};

% Bucket ticks
\foreach \t in {0,1,2,3} {
    \draw[black!40] (\t*\bx, \yb+0.08) -- (\t*\bx, \yb-0.08);
    \node[lbl, below=1pt, black!60] at (\t*\bx, \yb-0.08) {\t};
}

% Bucket boundary lines
\foreach \b in {0,1,2,3} {
    \draw[blue!30, thin, dashed]
        (\b*\bx, \yb-1.4) -- (\b*\bx, \yb+2.2);
}

% ── Activity i aggregated ────────────────────────────────
%    ES_agg = floor(2/5) = 0,  LS_agg = ceil(8/5) = 2,  d_agg = ceil(3/5) = 1
\fill[teal!50, rounded corners=2pt]
    (0*\bx, \yb+0.5) rectangle (1*\bx, \yb+1.0);
\node[font=\scriptsize, text=white, font=\bfseries]
    at (0.5*\bx, \yb+0.75) {$i$\; ($d_i^{\mathrm{agg}}{=}1$)};

% ES and LS markers
\draw[teal!70, thick] (0*\bx, \yb+0.35) -- (0*\bx, \yb+1.15);
\draw[teal!70, thick] (2*\bx, \yb+0.35) -- (2*\bx, \yb+1.15);
\node[lbl, teal!80, above=1pt] at (0*\bx, \yb+1.15)
    {$\lfloor 2/5 \rfloor{=}0$};
\node[lbl, teal!80, above=1pt] at (2*\bx, \yb+1.15)
    {$\lceil 8/5 \rceil{=}2$};

% Feasible window bracket
\draw[teal!40, <->, thick]
    (0*\bx, \yb+0.35) -- (2*\bx, \yb+0.35);

% ── Activity j aggregated ────────────────────────────────
%    ES_agg = floor(7/5) = 1,  LS_agg = ceil(12/5) = 3,  d_agg = ceil(4/5) = 1
\fill[violet!50, rounded corners=2pt]
    (1*\bx, \yb-1.1) rectangle (2*\bx, \yb-0.6);
\node[font=\scriptsize, text=white, font=\bfseries]
    at (1.5*\bx, \yb-0.85) {$j$\; ($d_j^{\mathrm{agg}}{=}1$)};

% ES and LS markers
\draw[violet!70, thick] (1*\bx, \yb-0.45) -- (1*\bx, \yb-1.25);
\draw[violet!70, thick] (3*\bx, \yb-0.45) -- (3*\bx, \yb-1.25);
\node[lbl, violet!80, below=1pt] at (1*\bx, \yb-1.25)
    {$\lfloor 7/5 \rfloor{=}1$};
\node[lbl, violet!80, below=1pt] at (3*\bx, \yb-1.25)
    {$\lceil 12/5 \rceil{=}3$};

% Feasible window bracket
\draw[violet!40, <->, thick]
    (1*\bx, \yb-0.45) -- (3*\bx, \yb-0.45);

\end{tikzpicture}
\caption{Time aggregation with factor $\alpha = 5$.
         Activities $i$ and $j$ are shown on the physical
         timeline~(a) and the aggregated timeline~(b).
         Coloured bars show each activity at its earliest start;
         double-headed arrows mark the feasible window $[ES, LS]$.
         Floor division maps $ES$ into buckets; ceiling division
         maps $LS$ and durations to preserve feasibility.}
\label{fig:time_agg}
\end{figure}


\subsubsection*{Time Window Transformation}

The earliest and latest start times are mapped from physical
periods to bucket indices:
\begin{align}
    ES_i^{\mathrm{agg}} &= \left\lfloor ES_i / \alpha
                           \right\rfloor
    \label{eq:es_agg} \\[4pt]
    LS_i^{\mathrm{agg}} &= \left\lceil LS_i / \alpha
                           \right\rceil
    \label{eq:ls_agg}
\end{align}
Floor division is used for $ES_i$ because a physical period
falling inside a bucket means the activity can begin in that
bucket. Ceiling division is used for $LS_i$ to avoid excluding
valid start positions: if $LS_i$ falls in the interior of a
bucket, rounding down would eliminate the last feasible start
period, potentially making the instance artificially infeasible.
The aggregated horizon is $T_{\mathrm{agg}} = \lceil T / \alpha
\rceil$.

\subsubsection*{Duration Transformation}

Activity durations are rounded up to ensure that the aggregated
duration covers at least as many buckets as the original activity
occupies:
\begin{equation}
    d_i^{\mathrm{agg}} = \left\lceil d_i / \alpha \right\rceil
    \label{eq:d_agg}
\end{equation}
Rounding down would undercount the number of buckets the activity
occupies, violating the work-balance constraint that requires the
activity to be active for enough periods to deliver its total
workload.

\subsubsection*{Resource Capacity and Rate Bound Transformation}

Each aggregated bucket spans $\alpha$ physical periods. The
per-period resource capacity $R_k$ therefore accumulates over the
bucket:
\begin{equation}
    R_k^{\mathrm{agg}} = \alpha \cdot R_k
    \label{eq:cap_agg}
\end{equation}
The upper rate bound scales identically,
$\overline{r}^{\,\mathrm{agg}} = \alpha \cdot \overline{r}$,
since the maximum resource consumption per bucket is the sum of
$\alpha$ per-period maxima.

The lower rate bound is set to zero rather than
$\alpha \cdot \underline{r}$:
\begin{equation}
    \underline{r}^{\,\mathrm{agg}} = 0
    \label{eq:lr_agg}
\end{equation}
This relaxation is necessary because scaling $\underline{r}$ by
$\alpha$ can render activities with small total workload
infeasible. Consider an activity with $\underline{r} = 1$ and
total work $W_{ik}^{\pi} = 3$: if $\alpha = 5$, the scaled lower
bound would be $\alpha \cdot \underline{r} = 5$, requiring the
activity to deliver at least 5 units of work per bucket, which is more
than its total work in a single bucket. Setting
$\underline{r}^{\,\mathrm{agg}} = 0$ ensures that the aggregated
problem is a \emph{relaxation} of the original: every schedule
feasible in physical time remains feasible in aggregated time,
so the search space of the aggregated problem includes all
solutions of the original problem.

\subsubsection*{Schedule Expansion}

After the R\&F loop produces a solution $S^{\mathrm{agg}}$ on
the aggregated grid, the schedule is expanded back to physical
time by mapping each aggregated start to the beginning of its
bucket:
\begin{equation}
    S_i^{\mathrm{phys}} = \alpha \cdot S_i^{\mathrm{agg}}
    \label{eq:expand}
\end{equation}
The expanded start times preserve precedence feasibility:
if $S_j^{\mathrm{agg}} + d_j^{\mathrm{agg}} \leq
S_i^{\mathrm{agg}}$ holds for every arc $(j, i) \in E$ in the
aggregated model, then $\alpha \cdot S_j^{\mathrm{agg}} +
d_j \leq \alpha \cdot S_j^{\mathrm{agg}} +
\alpha \cdot d_j^{\mathrm{agg}} \leq \alpha \cdot
S_i^{\mathrm{agg}}$, so the precedence constraint is satisfied
in physical time as well. Resource feasibility of the expanded
schedule is verified by the final LP pass described in
Section~\ref{subsec:final_verify}. The expanded start times
are aligned to bucket boundaries (multiples of $\alpha$), which
may be slightly conservative: the makespan of the expanded
schedule is bounded by
\begin{equation}
    C_{\max}^{\mathrm{phys}} \leq
    C_{\max}^{\mathrm{opt}} + \alpha \cdot |P^*|
    \label{eq:makespan_bound}
\end{equation}
where $|P^*|$ is the number of activities on the critical path
whose durations were rounded up during aggregation. The bound follows from the observation that each such activity
contributes at most $\alpha$ extra periods after expansion: its
aggregated duration satisfies
$\alpha \lceil d_i / \alpha \rceil \leq d_i + \alpha$, and the
critical-path makespan is the sum of durations along the longest
precedence chain. In practice, the inflation is small because
most critical-path activities have durations that are
near-multiples of $\alpha$.

\subsubsection*{Warm-Start Handling}

A physical warm-start schedule $S^{\mathrm{SSGS}}$ cannot simply
be translated to the aggregated grid via floor division
$S_i^{\mathrm{agg}} = \lfloor S_i^{\mathrm{SSGS}} / \alpha
\rfloor$, because rounding can create precedence violations on the
coarser grid. Consider two activities $j \to i$ with
$S_j^{\mathrm{SSGS}} = \alpha\tau - 1$ and
$d_j^{\mathrm{agg}} = 1$: in physical time, $j$ finishes at
$\alpha\tau$ and $i$ starts at $\alpha\tau$, which is a tight but valid
sequence. After floor division, both activities may map to the
same bucket $\tau - 1$, creating an apparent overlap. To avoid
this, the SSGS is re-run directly on the aggregated instance,
producing a warm-start that is guaranteed to be
precedence-feasible on the coarser grid.

Table~\ref{tab:time_agg_summary} summarises all transformations.

\begin{table}[ht]
\centering
\renewcommand{\arraystretch}{1.4}
\caption{Time aggregation transformations by factor $\alpha$}
\label{tab:time_agg_summary}
\begin{tabular}{lll}
\hline
\textbf{Quantity} &
\textbf{Aggregated value} &
\textbf{Rationale} \\
\hline
$ES_i$    & $\lfloor ES_i / \alpha \rfloor$
          & Activity can start in this bucket \\
$LS_i$    & $\lceil LS_i / \alpha \rceil$
          & Do not exclude valid start positions \\
$d_i$     & $\lceil d_i / \alpha \rceil$
          & Do not undercount occupied buckets \\
$T$       & $\lceil T / \alpha \rceil$
          & Horizon shrinks proportionally \\
$R_k$     & $\alpha \cdot R_k$
          & $\alpha$ periods of capacity per bucket \\
$\overline{r}$
          & $\alpha \cdot \overline{r}$
          & $\alpha$ periods of rate per bucket \\
$\underline{r}$
          & $0$
          & Prevent infeasibility from rounding \\
$w_{ik}^{\pi}$
          & unchanged
          & Total work is invariant \\
$(V, E)$  & unchanged
          & Precedence is time-independent \\
\hline
\end{tabular}
\end{table}

% ────────────────────────────────────────────────────────────
\subsection{Per-Group Scenario Reduction}
\label{subsec:scen_red}
% ────────────────────────────────────────────────────────────

The second scalability technique reduces the number of stochastic
scenarios used in each group subproblem. The SAA-FSRCPSP model
contains $|S|$ copies of the resource allocation variables
$r_{ikt\pi}$ and the activity indicator variables $y_{it\pi}$,
so model size scales linearly with~$|S|$. With $|S| = 10$
scenarios, these variables dominate: for $n = 60$ instances,
$y$ and $r$ together account for over $90\%$ of the total
variable count.

Not every scenario matters equally for every group. When the
R\&F loop processes group~$G_g$, only that group's activities
carry binary (free) variables; everything else is either fixed
or relaxed. Workload variation outside~$G_g$ cannot influence
the combinatorial choices in this subproblem, since the
corresponding variables are not free. It therefore makes sense
to perform a \emph{group-local} reduction: for each group, pick
a subset of $n_{\mathrm{keep}}$ representative scenarios based
solely on the workload variation of that group's activities.

\subsubsection*{Group-Local Workload Vectors}

For group $G_g$ and scenario $\pi$, the workload vector is the
concatenation of all resource demands for the group's activities:
\begin{equation}
    \mathbf{w}_{G_g}^{\pi} =
    \bigl(w_{ik}^{\pi}\bigr)_{i \in G_g,\; k \in K}
    \;\in \mathbb{R}^{|G_g| \cdot |K|}
    \label{eq:scen_vec}
\end{equation}
Only activities in~$G_g$ appear because the rest have fixed or
relaxed start times and do not interact with the
scenario-dependent constraints in the same way.

\subsubsection*{PAM K-Medoids Selection}

The $|S|$ workload vectors
$\{\mathbf{w}_{G_g}^{1}, \ldots, \mathbf{w}_{G_g}^{|S|}\}$ are
clustered into $n_{\mathrm{keep}}$ groups using the Partitioning
Around Medoids (PAM) algorithm \parencite{Kaufman1990medoids}.
Unlike K-Means \parencite{MacQueen1967}, PAM constrains cluster
centres to be \emph{actual data points} rather than synthetic
averages. This matters here because a K-Means centroid may not
correspond to any real scenario, and feeding an artificial
workload vector into the MIP could violate structural assumptions
of the FSRCPSP constraints; for instance, it may produce non-integer
demands where integer values are expected.

\paragraph{PAM objective.}
PAM seeks a medoid set
$M = \{m_1, \ldots, m_{n_{\mathrm{keep}}}\} \subseteq
\{1, \ldots, |S|\}$ that minimises the total assignment cost:
\begin{equation}
    \mathrm{TD}(M) =
    \sum_{\pi=1}^{|S|}
    \min_{c \in \{1,\ldots,n_{\mathrm{keep}}\}}
    \bigl\|\mathbf{w}_{G_g}^{\pi} -
           \mathbf{w}_{G_g}^{m_c}\bigr\|_2
    \label{eq:pam_objective}
\end{equation}
Because every candidate centre must be an element of the data
set, minimising equation~\eqref{eq:pam_objective} is a
combinatorial problem. PAM solves it heuristically in two
phases.

\paragraph{Build phase.}
The initial medoid set is constructed greedily. The first medoid
is the scenario that minimises the sum of distances to all
others:
\begin{equation}
    m_1 = \arg\min_{j \in \{1,\ldots,|S|\}}
    \sum_{\pi=1}^{|S|}
    \bigl\|\mathbf{w}_{G_g}^{j} -
           \mathbf{w}_{G_g}^{\pi}\bigr\|_2
    \label{eq:pam_build_first}
\end{equation}
Each subsequent medoid is chosen by a farthest-point heuristic:
the scenario whose minimum distance to any existing medoid is
largest, i.e.\
$m_{c+1} = \arg\max_{o \notin M} \min_{m \in M}
\|\mathbf{w}_{G_g}^{o} - \mathbf{w}_{G_g}^{m}\|_2$.
This ensures that the initial medoids are well-separated across
the workload space. After $n_{\mathrm{keep}}$ medoids have
been placed, every scenario is assigned to its closest medoid.

\paragraph{Swap phase.}
The algorithm then iterates over all pairs $(m_c, o)$ where
$m_c$ is a current medoid and $o$ is any non-medoid, and checks
whether replacing $m_c$ with $o$ decreases $\mathrm{TD}(M)$.
If so, the swap is performed and cluster assignments are
updated. The phase repeats until a full pass produces no
improving swap, at which point the algorithm has reached a local
minimum of~$\mathrm{TD}$.

\paragraph{Complexity.}
Each swap evaluation requires $O(|S|)$ distance lookups against
a precomputed $|S| \times |S|$ distance matrix, and there are
at most $O(|S| \cdot n_{\mathrm{keep}})$ candidate swaps per
iteration. Building the distance matrix costs $O(|S|^2 \cdot
|G_g| \cdot |K|)$. For $|S| \leq 10$, the entire procedure
finishes in microseconds, negligible next to the MIP solve.

\paragraph{Medoid output.}
Each selected medoid $\pi_c^*$ is a real scenario from the
original set:
\begin{equation}
    \pi_c^* = m_c \in \{1, \ldots, |S|\},
    \qquad c = 1, \ldots, n_{\mathrm{keep}}
    \label{eq:medoid}
\end{equation}
so all downstream constraints (work balance, rate bounds, chance
constraints) remain structurally valid. The reduced scenario set
for group $G_g$ is $S_g^{\mathrm{red}} = \{\pi_{1}^*, \ldots,
\pi_{n_{\mathrm{keep}}}^*\}$, and the subproblem MIP is rebuilt
using only these scenarios. The scenario indices are renumbered
$1, \ldots, n_{\mathrm{keep}}$ to maintain the contiguous
1-indexed numbering expected by the model builder. With
$n_{\mathrm{keep}} = 5$ and $|S| = 10$, the scenario dimension
of each subproblem is halved, reducing the $y$ and $r$ variable
counts by $50\%$.

\subsubsection*{Alternative: Heitsch--R\"omisch Backward Reduction}

An alternative to clustering-based selection is the backward
scenario reduction algorithm of \textcite{Heitsch2003}. This
method greedily deletes scenarios one at a time, each time
removing the scenario whose deletion incurs the smallest
redistribution cost:
\begin{equation}
    \pi^* = \arg\min_{\pi \in S_{\mathrm{active}}}
    \; p_{\pi} \cdot
    \min_{j \in S_{\mathrm{active}} \setminus \{\pi\}}
    \bigl\|\mathbf{w}^{\pi} - \mathbf{w}^{j}\bigr\|_2
    \label{eq:backward_deletion}
\end{equation}
where $p_{\pi}$ is the probability weight of scenario~$\pi$
(uniform $1/|S|$ in this work). After deletion, the weight
$p_{\pi^*}$ is redistributed to the nearest remaining scenario.
The procedure repeats until $n_{\mathrm{keep}}$ scenarios
remain. Under uniform weights, the retained set minimises
$\sum_{\pi \in \text{deleted}} \min_{j \in \text{retained}}
\|\mathbf{w}^{\pi} - \mathbf{w}^{j}\|$, which bounds the
type-1 Wasserstein distance between the original and reduced
distributions \parencite{Heitsch2003}.

\paragraph{Comparison with PAM.}
The backward method optimises a different objective
(Wasserstein distance) from PAM (total assignment cost), but
for $|S| \leq 10$ both produce nearly identical retained sets.
The experiments in Chapter~\ref{chap:experiments},
Section~\ref{sec:scen_red_exp} confirm this: on j120 instances
with global reduction, PAM achieves a mean gap of~$31.83\%$
and the backward method~$32.01\%$, a difference well within
instance-level noise. Since both algorithms produce comparable results at this scale,
PAM is used as the default: it is the simpler of the two to
implement and the more widely used algorithm in the clustering
literature, making the implementation easier to reproduce.

A per-group variant of the backward method (applying
equation~\eqref{eq:backward_deletion} to the group-local
vectors $\mathbf{w}_{G_g}^{\pi}$ instead of the full-instance
vectors) would be straightforward to implement and is expected
to yield results close to per-group PAM, since the global
comparison already shows the two selection algorithms perform
comparably. The experiments therefore focus on the global-vs.\
per-group distinction, which is the dominant factor: per-group
K-medoids achieves a mean gap of~$10.30\%$, compared to
$31$--$33\%$ for all four global methods
(Section~\ref{sec:scen_red_exp}).

\subsubsection*{Why Per-Group Selection Matters}

Because the workload vectors are built from group-local
activities, different groups typically select different
representative scenarios. A group whose activities show high
workload variation on resource~$k = 1$ may pick scenarios that
capture extremes on that resource, while a group dominated by
resource~$k = 2$ variation may pick an entirely different subset.
Each subproblem therefore sees the scenarios most relevant to its
own combinatorial decisions, rather than a single global subset
that may be a poor fit for any particular group.
Figure~\ref{fig:pergroup_scen} illustrates this: the same
10~scenarios yield different retained subsets for two groups
with different resource-demand profiles.

% ============================================================
% TikZ diagram: Per-Group Scenario Reduction
% Shows 10 scenarios reduced to different subsets for two groups
%
% Include in Chapter 5 with:
%   % ============================================================
% TikZ diagram: Per-Group Scenario Reduction
% Shows 10 scenarios reduced to different subsets for two groups
%
% Include in Chapter 5 with:
%   % ============================================================
% TikZ diagram: Per-Group Scenario Reduction
% Shows 10 scenarios reduced to different subsets for two groups
%
% Include in Chapter 5 with:
%   \input{figures/pergroup_scenario_reduction.tex}
%
% Required in preamble:
%   \usepackage{tikz}
%   \usetikzlibrary{arrows.meta, positioning, calc}
% ============================================================
\begin{figure}[ht]
\centering
\begin{tikzpicture}[
    >=Stealth,
    lbl/.style={font=\scriptsize},
    scen/.style={circle, minimum size=0.55cm,
                 draw=black!30, inner sep=0pt,
                 font=\scriptsize\bfseries},
    kept/.style={scen, fill=#1, text=white, draw=#1!70,
                 line width=1pt},
    dropped/.style={scen, fill=black!6, text=black!35,
                    draw=black!15},
    groupbox/.style={rectangle, rounded corners=6pt,
                     draw=#1!50, fill=#1!5,
                     minimum width=3.2cm, minimum height=1.0cm,
                     align=center, font=\small\bfseries,
                     text=#1!80},
    arrpam/.style={->, thick, draw=black!40,
                   shorten >=3pt, shorten <=3pt},
]

% ================================================================
%  Full scenario set (top, centred)
% ================================================================
\node[font=\small\bfseries, black!70] at (5, 5.2)
    {Full scenario set $S = \{\pi_1, \ldots, \pi_{10}\}$};

% 10 scenario circles in a row
\foreach \s in {1,...,10} {
    \node[scen, fill=black!8] (S\s) at (\s, 4.3) {\s};
}

% ================================================================
%  Group G_1 (left) — resource k=1 dominant
% ================================================================
\node[groupbox=teal] (G1) at (3, 2.4)
    {Group $G_1$\\[-2pt]
     {\scriptsize resource $k{=}1$ dominant}};

% Arrow from full set to G1
\draw[arrpam] (3, 3.9) -- (G1.north)
    node[midway, left=4pt, font=\scriptsize, black!50,
         align=right] {PAM on\\$\mathbf{w}_{G_1}^{\pi}$};

% Reduced set for G1: keep scenarios 2, 5, 8 (example)
\node[font=\scriptsize, black!60, anchor=south]
    at (3, 0.95) {$S_1^{\mathrm{red}} = \{2, 5, 8\}$};

\foreach \s/\x in {1/0.5, 2/1.5, 3/2.5, 4/3.5, 5/4.5} {
    \pgfmathtruncatemacro{\sid}{
        ifthenelse(\s==1, 1,
        ifthenelse(\s==2, 2,
        ifthenelse(\s==3, 5,
        ifthenelse(\s==4, 8, 10))))}
    % manual placement
}

% Manual placement — cleaner
\node[dropped]       (G1S1)  at (1.0, 0.5) {1};
\node[kept=teal!70]  (G1S2)  at (1.8, 0.5) {2};
\node[dropped]       (G1S3)  at (2.6, 0.5) {3};
\node[dropped]       (G1S4)  at (3.4, 0.5) {4};
\node[kept=teal!70]  (G1S5)  at (4.2, 0.5) {5};
\node[dropped]       (G1S6)  at (1.4, -0.3) {6};
\node[dropped]       (G1S7)  at (2.2, -0.3) {7};
\node[kept=teal!70]  (G1S8)  at (3.0, -0.3) {8};
\node[dropped]       (G1S9)  at (3.8, -0.3) {9};
\node[dropped]       (G1S10) at (4.6, -0.3) {10};

% ================================================================
%  Group G_2 (right) — resource k=2 dominant
% ================================================================
\node[groupbox=violet] (G2) at (8, 2.4)
    {Group $G_2$\\[-2pt]
     {\scriptsize resource $k{=}2$ dominant}};

% Arrow from full set to G2
\draw[arrpam] (8, 3.9) -- (G2.north)
    node[midway, right=4pt, font=\scriptsize, black!50,
         align=left] {PAM on\\$\mathbf{w}_{G_2}^{\pi}$};

% Reduced set for G2: keep scenarios 1, 4, 9 (different subset!)
\node[font=\scriptsize, black!60, anchor=south]
    at (8, 0.95) {$S_2^{\mathrm{red}} = \{1, 4, 9\}$};

\node[kept=violet!70] (G2S1)  at (6.0, 0.5) {1};
\node[dropped]         (G2S2)  at (6.8, 0.5) {2};
\node[dropped]         (G2S3)  at (7.6, 0.5) {3};
\node[kept=violet!70]  (G2S4)  at (8.4, 0.5) {4};
\node[dropped]         (G2S5)  at (9.2, 0.5) {5};
\node[dropped]         (G2S6)  at (6.4, -0.3) {6};
\node[dropped]         (G2S7)  at (7.2, -0.3) {7};
\node[dropped]         (G2S8)  at (8.0, -0.3) {8};
\node[kept=violet!70]  (G2S9)  at (8.8, -0.3) {9};
\node[dropped]         (G2S10) at (9.6, -0.3) {10};

% ================================================================
%  Annotation: key insight
% ================================================================
\node[font=\scriptsize, black!50, align=center,
      text width=5cm] at (5.5, -1.1)
    {Different groups select different representative scenarios\\
     based on their own workload variation};

\end{tikzpicture}
\caption{Per-group scenario reduction. The same 10~scenarios are
         reduced to $n_{\mathrm{keep}}{=}3$ for each group, but
         PAM selects different medoids because the group-local
         workload vectors $\mathbf{w}_{G_g}^{\pi}$ differ.
         Coloured circles are retained; grey circles are dropped.}
\label{fig:pergroup_scen}
\end{figure}

%
% Required in preamble:
%   \usepackage{tikz}
%   \usetikzlibrary{arrows.meta, positioning, calc}
% ============================================================
\begin{figure}[ht]
\centering
\begin{tikzpicture}[
    >=Stealth,
    lbl/.style={font=\scriptsize},
    scen/.style={circle, minimum size=0.55cm,
                 draw=black!30, inner sep=0pt,
                 font=\scriptsize\bfseries},
    kept/.style={scen, fill=#1, text=white, draw=#1!70,
                 line width=1pt},
    dropped/.style={scen, fill=black!6, text=black!35,
                    draw=black!15},
    groupbox/.style={rectangle, rounded corners=6pt,
                     draw=#1!50, fill=#1!5,
                     minimum width=3.2cm, minimum height=1.0cm,
                     align=center, font=\small\bfseries,
                     text=#1!80},
    arrpam/.style={->, thick, draw=black!40,
                   shorten >=3pt, shorten <=3pt},
]

% ================================================================
%  Full scenario set (top, centred)
% ================================================================
\node[font=\small\bfseries, black!70] at (5, 5.2)
    {Full scenario set $S = \{\pi_1, \ldots, \pi_{10}\}$};

% 10 scenario circles in a row
\foreach \s in {1,...,10} {
    \node[scen, fill=black!8] (S\s) at (\s, 4.3) {\s};
}

% ================================================================
%  Group G_1 (left) — resource k=1 dominant
% ================================================================
\node[groupbox=teal] (G1) at (3, 2.4)
    {Group $G_1$\\[-2pt]
     {\scriptsize resource $k{=}1$ dominant}};

% Arrow from full set to G1
\draw[arrpam] (3, 3.9) -- (G1.north)
    node[midway, left=4pt, font=\scriptsize, black!50,
         align=right] {PAM on\\$\mathbf{w}_{G_1}^{\pi}$};

% Reduced set for G1: keep scenarios 2, 5, 8 (example)
\node[font=\scriptsize, black!60, anchor=south]
    at (3, 0.95) {$S_1^{\mathrm{red}} = \{2, 5, 8\}$};

\foreach \s/\x in {1/0.5, 2/1.5, 3/2.5, 4/3.5, 5/4.5} {
    \pgfmathtruncatemacro{\sid}{
        ifthenelse(\s==1, 1,
        ifthenelse(\s==2, 2,
        ifthenelse(\s==3, 5,
        ifthenelse(\s==4, 8, 10))))}
    % manual placement
}

% Manual placement — cleaner
\node[dropped]       (G1S1)  at (1.0, 0.5) {1};
\node[kept=teal!70]  (G1S2)  at (1.8, 0.5) {2};
\node[dropped]       (G1S3)  at (2.6, 0.5) {3};
\node[dropped]       (G1S4)  at (3.4, 0.5) {4};
\node[kept=teal!70]  (G1S5)  at (4.2, 0.5) {5};
\node[dropped]       (G1S6)  at (1.4, -0.3) {6};
\node[dropped]       (G1S7)  at (2.2, -0.3) {7};
\node[kept=teal!70]  (G1S8)  at (3.0, -0.3) {8};
\node[dropped]       (G1S9)  at (3.8, -0.3) {9};
\node[dropped]       (G1S10) at (4.6, -0.3) {10};

% ================================================================
%  Group G_2 (right) — resource k=2 dominant
% ================================================================
\node[groupbox=violet] (G2) at (8, 2.4)
    {Group $G_2$\\[-2pt]
     {\scriptsize resource $k{=}2$ dominant}};

% Arrow from full set to G2
\draw[arrpam] (8, 3.9) -- (G2.north)
    node[midway, right=4pt, font=\scriptsize, black!50,
         align=left] {PAM on\\$\mathbf{w}_{G_2}^{\pi}$};

% Reduced set for G2: keep scenarios 1, 4, 9 (different subset!)
\node[font=\scriptsize, black!60, anchor=south]
    at (8, 0.95) {$S_2^{\mathrm{red}} = \{1, 4, 9\}$};

\node[kept=violet!70] (G2S1)  at (6.0, 0.5) {1};
\node[dropped]         (G2S2)  at (6.8, 0.5) {2};
\node[dropped]         (G2S3)  at (7.6, 0.5) {3};
\node[kept=violet!70]  (G2S4)  at (8.4, 0.5) {4};
\node[dropped]         (G2S5)  at (9.2, 0.5) {5};
\node[dropped]         (G2S6)  at (6.4, -0.3) {6};
\node[dropped]         (G2S7)  at (7.2, -0.3) {7};
\node[dropped]         (G2S8)  at (8.0, -0.3) {8};
\node[kept=violet!70]  (G2S9)  at (8.8, -0.3) {9};
\node[dropped]         (G2S10) at (9.6, -0.3) {10};

% ================================================================
%  Annotation: key insight
% ================================================================
\node[font=\scriptsize, black!50, align=center,
      text width=5cm] at (5.5, -1.1)
    {Different groups select different representative scenarios\\
     based on their own workload variation};

\end{tikzpicture}
\caption{Per-group scenario reduction. The same 10~scenarios are
         reduced to $n_{\mathrm{keep}}{=}3$ for each group, but
         PAM selects different medoids because the group-local
         workload vectors $\mathbf{w}_{G_g}^{\pi}$ differ.
         Coloured circles are retained; grey circles are dropped.}
\label{fig:pergroup_scen}
\end{figure}

%
% Required in preamble:
%   \usepackage{tikz}
%   \usetikzlibrary{arrows.meta, positioning, calc}
% ============================================================
\begin{figure}[ht]
\centering
\begin{tikzpicture}[
    >=Stealth,
    lbl/.style={font=\scriptsize},
    scen/.style={circle, minimum size=0.55cm,
                 draw=black!30, inner sep=0pt,
                 font=\scriptsize\bfseries},
    kept/.style={scen, fill=#1, text=white, draw=#1!70,
                 line width=1pt},
    dropped/.style={scen, fill=black!6, text=black!35,
                    draw=black!15},
    groupbox/.style={rectangle, rounded corners=6pt,
                     draw=#1!50, fill=#1!5,
                     minimum width=3.2cm, minimum height=1.0cm,
                     align=center, font=\small\bfseries,
                     text=#1!80},
    arrpam/.style={->, thick, draw=black!40,
                   shorten >=3pt, shorten <=3pt},
]

% ================================================================
%  Full scenario set (top, centred)
% ================================================================
\node[font=\small\bfseries, black!70] at (5, 5.2)
    {Full scenario set $S = \{\pi_1, \ldots, \pi_{10}\}$};

% 10 scenario circles in a row
\foreach \s in {1,...,10} {
    \node[scen, fill=black!8] (S\s) at (\s, 4.3) {\s};
}

% ================================================================
%  Group G_1 (left) — resource k=1 dominant
% ================================================================
\node[groupbox=teal] (G1) at (3, 2.4)
    {Group $G_1$\\[-2pt]
     {\scriptsize resource $k{=}1$ dominant}};

% Arrow from full set to G1
\draw[arrpam] (3, 3.9) -- (G1.north)
    node[midway, left=4pt, font=\scriptsize, black!50,
         align=right] {PAM on\\$\mathbf{w}_{G_1}^{\pi}$};

% Reduced set for G1: keep scenarios 2, 5, 8 (example)
\node[font=\scriptsize, black!60, anchor=south]
    at (3, 0.95) {$S_1^{\mathrm{red}} = \{2, 5, 8\}$};

\foreach \s/\x in {1/0.5, 2/1.5, 3/2.5, 4/3.5, 5/4.5} {
    \pgfmathtruncatemacro{\sid}{
        ifthenelse(\s==1, 1,
        ifthenelse(\s==2, 2,
        ifthenelse(\s==3, 5,
        ifthenelse(\s==4, 8, 10))))}
    % manual placement
}

% Manual placement — cleaner
\node[dropped]       (G1S1)  at (1.0, 0.5) {1};
\node[kept=teal!70]  (G1S2)  at (1.8, 0.5) {2};
\node[dropped]       (G1S3)  at (2.6, 0.5) {3};
\node[dropped]       (G1S4)  at (3.4, 0.5) {4};
\node[kept=teal!70]  (G1S5)  at (4.2, 0.5) {5};
\node[dropped]       (G1S6)  at (1.4, -0.3) {6};
\node[dropped]       (G1S7)  at (2.2, -0.3) {7};
\node[kept=teal!70]  (G1S8)  at (3.0, -0.3) {8};
\node[dropped]       (G1S9)  at (3.8, -0.3) {9};
\node[dropped]       (G1S10) at (4.6, -0.3) {10};

% ================================================================
%  Group G_2 (right) — resource k=2 dominant
% ================================================================
\node[groupbox=violet] (G2) at (8, 2.4)
    {Group $G_2$\\[-2pt]
     {\scriptsize resource $k{=}2$ dominant}};

% Arrow from full set to G2
\draw[arrpam] (8, 3.9) -- (G2.north)
    node[midway, right=4pt, font=\scriptsize, black!50,
         align=left] {PAM on\\$\mathbf{w}_{G_2}^{\pi}$};

% Reduced set for G2: keep scenarios 1, 4, 9 (different subset!)
\node[font=\scriptsize, black!60, anchor=south]
    at (8, 0.95) {$S_2^{\mathrm{red}} = \{1, 4, 9\}$};

\node[kept=violet!70] (G2S1)  at (6.0, 0.5) {1};
\node[dropped]         (G2S2)  at (6.8, 0.5) {2};
\node[dropped]         (G2S3)  at (7.6, 0.5) {3};
\node[kept=violet!70]  (G2S4)  at (8.4, 0.5) {4};
\node[dropped]         (G2S5)  at (9.2, 0.5) {5};
\node[dropped]         (G2S6)  at (6.4, -0.3) {6};
\node[dropped]         (G2S7)  at (7.2, -0.3) {7};
\node[dropped]         (G2S8)  at (8.0, -0.3) {8};
\node[kept=violet!70]  (G2S9)  at (8.8, -0.3) {9};
\node[dropped]         (G2S10) at (9.6, -0.3) {10};

% ================================================================
%  Annotation: key insight
% ================================================================
\node[font=\scriptsize, black!50, align=center,
      text width=5cm] at (5.5, -1.1)
    {Different groups select different representative scenarios\\
     based on their own workload variation};

\end{tikzpicture}
\caption{Per-group scenario reduction. The same 10~scenarios are
         reduced to $n_{\mathrm{keep}}{=}3$ for each group, but
         PAM selects different medoids because the group-local
         workload vectors $\mathbf{w}_{G_g}^{\pi}$ differ.
         Coloured circles are retained; grey circles are dropped.}
\label{fig:pergroup_scen}
\end{figure}


% ────────────────────────────────────────────────────────────
\subsection{Final Verification Pass}
\label{subsec:final_verify}
% ────────────────────────────────────────────────────────────

Both scalability techniques introduce approximations: time
aggregation solves on a coarser grid whose resource constraints
are weaker than the physical ones, and per-group scenario
reduction solves each subproblem on a subset of
$n_{\mathrm{keep}} < |S|$ scenarios rather than the full set.
While the medoid selection ensures that the chosen scenarios are
representative, there is no formal guarantee that the resulting
schedule satisfies all resource and chance constraints on the
original, full-resolution instance. A final verification pass
closes this gap.

After the R\&F loop terminates, the algorithm rebuilds the full
SAA-FSRCPSP model using all $|S|$ original scenarios, fixes
\emph{all} activity start times to the schedule $S^*$ obtained
from the R\&F loop (i.e.\ $x^B_{i, S^*_i} = 1$ and
$x^B_{it} = 0$ for all $t \neq S^*_i$) and solves the resulting
resource allocation problem. With all binary variables pinned,
the remaining model is a continuous LP over the resource
allocation variables $r_{ikt\pi}$ and activity indicators
$y_{it\pi}$. If this LP is feasible and the integrality condition
$x^{\mathrm{eff}}_{it} \in \{0, 1\}$ is satisfied for all
activities, the schedule is certified as feasible on the full
scenario set. Otherwise, the algorithm reports the schedule as
not certified.

The overall flow (\emph{reduce}, \emph{solve}, \emph{verify}) is
the standard matheuristic pattern of solving an approximate model
first and checking the result on the original afterwards
\parencite{Maniezzo2010}. The verification adds little cost:
with all start times pinned, only a continuous LP remains, and
its difficulty is independent of the combinatorial scheduling
problem.

\subsubsection*{Combined Size Reduction}

The two techniques compose multiplicatively along independent
dimensions of the model. Time aggregation shrinks the
$t$-dimension by a factor of~$\alpha$ (from $T$ to
$\lceil T/\alpha \rceil$), and scenario reduction shrinks the
$\pi$-dimension by a factor of
$|S|/n_{\mathrm{keep}}$ (from $|S|$ to $n_{\mathrm{keep}}$).
For $\alpha = 5$, $|S| = 10$, and $n_{\mathrm{keep}} = 5$,
the combined reduction factor is
$5 \times 2 = 10$: each subproblem MIP is approximately
one-tenth the size of the original. This reduction is sufficient to bring
$n = 60$ instances within practical solve times while preserving
solution quality through the final verification pass.

The extended procedure incorporating both scalability techniques
is given in Algorithm~\ref{alg:group_rf_extended}.

% ── Algorithm 9 ──────────────────────────────────────────────
\begin{algorithm}[H]
\caption{Group-based R\&F with time aggregation and
         per-group scenario reduction}
\label{alg:group_rf_extended}
\begin{algorithmic}[1]
\State Determine greedy solution $S$ via serial SGS
\State Compute $\tilde{X}$ using Algorithm~\ref{alg:features}
\State Compute $\mathcal{G}=[G_1,\ldots,G_m]$
       using Algorithm~\ref{alg:grouping}
\If{$\alpha > 1$}
    \State \parbox[t]{0.80\linewidth}{%
           Apply time aggregation
           (Table~\ref{tab:time_agg_summary}) to obtain
           aggregated instance with
           $T_{\mathrm{agg}} = \lceil T/\alpha \rceil$}
    \State Re-run SSGS on aggregated instance to obtain
           feasible warm-start $S^{\mathrm{agg}}$
\EndIf
\For{$g=1$ \textbf{to} $|\mathcal{G}|$}
    \If{per-group scenario reduction enabled}
        \State \parbox[t]{0.78\linewidth}{%
               Build group-local workload vectors
               $\mathbf{w}_{G_g}^{\pi}$ for all $\pi \in S$
               using only activities in $G_g$
               (equation~\ref{eq:scen_vec})}
        \State \parbox[t]{0.78\linewidth}{%
               Select $n_{\mathrm{keep}}$ representative
               scenarios via PAM K-medoids
               (equation~\ref{eq:medoid})}
        \State Rebuild MIP using reduced scenario set
               $S_g^{\mathrm{red}}$
    \EndIf
    \For{$i\in G_g$}
        \State $w_i \leftarrow \max(2,\,\lfloor \alpha_w \cdot
               (LS_i-ES_i)\rfloor)$
        \State $\widetilde{W}_i \leftarrow
               [\max\{ES_i,\,S'_i-w_i\},\;
                \min\{LS_i,\,S'_i+w_i\}]$
    \EndFor
    \State Solve subproblem with narrow windows
           ($\alpha_w = 0.15$)
    \If{no feasible solution found}
        \State Set $\widetilde{W}_i=[ES_i,\,LS_i]$
               for all $i\in G_g$; re-solve with full windows
    \EndIf
    \If{feasible solution $S'$ found}
        \State Fix $x^B_{i,S'_i}:=1$ for all $i \in G_g$;\;
               update $S \leftarrow S'$
        \If{\texttt{stop\_on\_first\_feasible}}
            \State \textbf{break}
        \EndIf
    \Else
        \State Fix $i\in G_g$ at current $S'_i$
               \Comment{fallback}
    \EndIf
\EndFor
\If{$\alpha > 1$}
    \State $S_i^{\mathrm{phys}} \leftarrow \alpha \cdot
           S_i^{\mathrm{agg}}$ for all $i$
           \Comment{expand to physical time}
\EndIf
\If{time aggregation or scenario reduction was used}
    \State \parbox[t]{0.78\linewidth}{%
           Rebuild MIP on the original (non-aggregated) instance
           with all $|S|$ scenarios;
           fix all $x^B$ to $S^*$; solve resource allocation
           LP to verify feasibility on full instance}
\EndIf
\State \Return Solution $S$
\end{algorithmic}
\end{algorithm}

% ------------------------------------------------------------
\section{Complexity Analysis}
\label{sec:grouprf_complexity}
% ------------------------------------------------------------
The group-based approach reduces the number of MIP subproblems
from $n+2$ (baseline) to $m = |\mathcal{G}|$ (group-based), where
$m$ depends on the instance structure and the grouping parameters.
For the benchmark instances used in this thesis
($n \in \{10, 30, 60\}$, $|K| \in \{1, 2, 3, 4\}$), the grouping
pipeline produces between 2 and 12 groups, yielding a reduction of
$60\%$--$80\%$ in the number of subproblems relative to the
baseline.

Each individual subproblem is larger than in the baseline (it
contains $|G_g|$ binary activities instead of one), but this cost
is offset by tighter root-node relaxations that accelerate
branch-and-bound convergence \parencite{pochet2006}.

The grouping overhead is negligible compared with MIP solve
times. Feature extraction runs in
$O(n^2 + n \cdot |K| \cdot |S|)$, Ward clustering in
$O(n^2 \log n)$ per precedence level \parencite{Murtagh2012},
and the merge safety check in
$O(|V| \cdot (|V| + |E|))$. For the instance sizes considered
here, the entire pipeline completes in under one second.

Termination criteria further reduce the number of subproblems
solved in practice. With \texttt{stop\_on\_first\_feasible}
enabled, the loop halts as soon as the first group subproblem
returns a feasible integer solution; when disabled, the
\texttt{max\_no\_improve} counter stops after a configurable
number of consecutive non-improving groups. In practice, the
algorithm typically solves one to three subproblems.

For large instances ($n = 60$), the scalability extensions of
Section~\ref{sec:scalability} provide additional reductions along
the time and scenario dimensions of the model. With time
aggregation factor $\alpha$ and per-group scenario reduction from
$|S|$ to $n_{\mathrm{keep}}$ scenarios, each subproblem shrinks
by a factor of approximately
$\alpha \cdot |S|/n_{\mathrm{keep}}$ relative to the
non-aggregated model. For the configuration used in the $n = 60$
experiments ($\alpha = 5$, $n_{\mathrm{keep}} = 5$, $|S| = 10$),
this yields a $10\times$ reduction in per-subproblem MIP size,
bringing instances that would otherwise require hours of
computation into the range of 5--15 minutes.

\bigskip
\noindent
In summary, the group-based R\&F pipeline replaces the
one-activity-at-a-time baseline with a principled decomposition
driven by feature similarity: Ward clustering partitions
activities into precedence-safe groups, and the R\&F loop solves
one group per subproblem with slack-proportional windows and
early termination. For instances beyond $n = 30$, time
aggregation and per-group scenario reduction keep the
subproblems tractable without sacrificing solution quality.
Chapter~\ref{chap:experiments} evaluates this pipeline against
the baseline across four instance sizes
($n \in \{10, 30, 60, 120\}$) and four resource configurations
($|K| \in \{1, 2, 3, 4\}$).
